% !TeX encoding = UTF-8
% Ce fichier contient le code de l'extension "chemfig"
%
%%%%%%%%%%%%%%%%%%%%%%%%%%%%%%%%%%%%%%%%%%%%%%%%%%%%%%%%%%%%%%%%%%%%%%
%                                                                    %
\def\CFname                    {chemfig}                             %
\def\CFver                       {1.81}                              %
%                                                                    %
\def\CFdate                   {2026/09/01}                           %
%                                                                    %
%%%%%%%%%%%%%%%%%%%%%%%%%%%%%%%%%%%%%%%%%%%%%%%%%%%%%%%%%%%%%%%%%%%%%%
%
%____________________________________________________________________
% Author     : Christian Tellechea                                   |
% Status     : Maintained                                            |
% Email      : unbonpetit@netc.fr                                    |
% Package URL: https://www.ctan.org/pkg/chemfig                      |
% Copyright  : Christian Tellechea 2010-2026                         |
% Licence    : Released under the LaTeX Project Public License v1.3c |
%              or later, see http://www.latex-project.org/lppl.txt   |
% Files      : 1) chemfig.tex                                        |
%              2) chemfig.sty                                        |
%              3) README                                             |
%              4) chemfig_doc_fr.tex                                 |
%              5) chemfig_doc_fr.pdf                                 |
%              6) chemfig_doc_en.tex                                 |
%              7) chemfig_doc_fr.pdf                                 |
%              8) chemfig-lewis.tex                                  |
%              9) changelog.tex                                      |
%--------------------------------------------------------------------

%%%%%%%%%%%%%%%%%%%%%%%%%%%%%%%%%%%%%%%%%%%%%%%%%%%%%%%%%%%%%%%%%%%%%%
%%%                       P R É A L A B L E                        %%%
%%%%%%%%%%%%%%%%%%%%%%%%%%%%%%%%%%%%%%%%%%%%%%%%%%%%%%%%%%%%%%%%%%%%%%
\csname CFloadonce\endcsname
\let\CFloadonce\endinput
%-------------------------- Annonce package --------------------------
\ifdefined\CFfromsty\else
	\immediate\write -1 {%
	Package: \CFname\space\CFdate\space\space v\CFver\space\space
	Draw molecule with an easy syntax (CT)}%
\fi

%-------------------------- Régime catcodes --------------------------
\begingroup
	\def\X#1{\catcode\number`#1=\number\catcode`#1\relax}
	\expandafter\xdef\csname CF_restorecatcode\endcsname{\X\[\X\]\X\:\X\(\X\)\X\,\X\-\X\=\X\~\X\!\X\?\X\<\X\>\X\;\X\*\X\|\X\@\X\ \X\_}
\endgroup
\catcode`\_11
\def\CF_setcatcodes#1{%
	\CF_setcatcodes_a#1,00,% bugfix 1.81
}
\def\CF_setcatcodes_a#1#2,{%
	\ifnum#2>0 % bugfix 1.81
		\catcode`#1#2\relax
		\expandafter\CF_setcatcodes_a
	\fi
}
\CF_setcatcodes{\[12,\]12,\:12,\(12,\)12,\,12,\-12,\=12,\~12,\!12,\?12,\<12,\>12,\;12,\*12,\|12,\@12,\#6,\ 10}

%-------------------- Vérification des prérequis ---------------------
\def\CF_error#1{\errmessage{Package \CFname\space Error: #1.}}
\def\CF_warning#1{\immediate\write-1{Package \CFname\space Warning: #1^^J}}
\def\CF_checkprimitive#1#2{% Vérifie que #1 est une primitive et sinon, émet le message #2 et exécute \endinput
	\begingroup
		\edef\__tempa{\meaning#1}%
		\edef\__tempb{\string#1}%
		\expandafter
	\endgroup
	\ifx\__tempa\__tempb\else
		\CF_error{#2}%
		\expandafter\expandafter\endinput
	\fi
}
\CF_checkprimitive\eTeXversion{You are not using an eTeX engine, \CFname\space cannot work.}
\CF_checkprimitive\expanded{the \string\expanded\space primitive is not provided by your TeX engine, \CFname\space cannot work.}

%------------------------ Chargement simplekv ------------------------
\input simplekv.tex

%-------------------------- Chargement tikz --------------------------
\unless\ifdefined\tikzpicture
	\begingroup\def\CF_temp{\endgroup\input tikz.tex\relax}%
	\expandafter\CF_temp
\fi
\usetikzlibrary{arrows.meta}

%--------------------------- Allocations -----------------------------
\newcount\CF_cnt_atomgroup
\newcount\CF_cnt_group
\newcount\CF_cnt_atom
\newcount\CF_cnt_cycle
\newcount\CF_cnt_cyclebonds
\newcount\CF_cnt_compound
\newcount\CF_cnt_save
\newcount\CF_cnt_hreac

\newif\ifCF_in_cycle
\newif\ifCF_cycle_arc
\newif\ifCF_define_submol
\newif\ifCF_adjust_name_dp
\newif\ifCF_macro_fixed_bond_length
\newif\ifCF_compound_is_chemfig

\newdimen\CF_dim_tmp
\newdimen\CF_dim_turnangle
\newdimen\CF_dim_arrowsize
\newdimen\CF_zero
\CF_zero=0pt

\newbox\CF_box
\newbox\CF_box_name
\newbox\CF_box_zerodim
\newbox\CF_box_charge
\newbox\CF_box_harrowup
\newbox\CF_box_harrow_down

\newtoks\CF_toks_subst

\newwrite\CF_write

%-------------------------- Petites macros ---------------------------
\let\CF_begintikzpicture\tikzpicture
\let\CF_endtikzpicture  \endtikzpicture
\def\CF_quark{\CF_quark}
\def\CF_execfirst#1#2{#1}
\def\CF_execsecond#1#2{#2}
\def\CF_id#1{#1}
\def\CF_gobarg#1{}
\def\CF_gobtwoargs#1#2{}
\def\CF_firsttonil#1#2\_nil{#1}
\def\CF_sanitizelastitem#1,\empty#2\_nil{#1}
\def\CF_gobtikzinstruction#1;{}
\def\CF_makeother#1{\catcode`#1=12\relax}
\def\CF_lettoken#1#2{\let#1= #2}
\CF_lettoken\CF_sptoken{ }
\def\CF_test_ok\fi\CF_execsecond#1#2{\fi#1}
\def\CF_ifx#1#2{\ifx#1#2\CF_test_ok\fi\CF_execsecond}
\def\CF_ifempty#1{\ifx\empty#1\empty\CF_test_ok\fi\CF_execsecond}
\def\CF_ifnum#1{\ifnum#1\CF_test_ok\fi\CF_execsecond}
\def\CF_ifinsidetikz{\ifdefined\pgfpictureid\CF_test_ok\fi\CF_execsecond}
\def\CF_ifzerodim#1{%r bugfix 1.71 : rollback à v1.66 !
	\setbox\CF_box_zerodim\hbox{\pgfinterruptpicture\printatom{#1}\endpgfinterruptpicture}% bugfix 1.53
	\CF_ifnum{1\ifdim\wd\CF_box_zerodim=\CF_zero0\fi\ifdim\ht\CF_box_zerodim=\CF_zero0\fi\ifdim\dp\CF_box_zerodim=\CF_zero0\fi=1000 }
}
\def\CF_doifempty#1{\ifx\empty#1\empty\expandafter\CF_id\else\expandafter\CF_gobarg\fi}
\def\CF_doifnotempty#1{\ifx\empty#1\empty\expandafter\CF_gobarg\else\expandafter\CF_id\fi}
\def\CF_gobtonil#1\_nil{}
\def\CF_stripsp#1{%
	\long\def\CF_stripsp##1##2{\expanded{\CF_stripsp_i\_marksp##1\__nil\_marksp#1\_marksp\_nil{##2}}}%
	\long\def\CF_stripsp_i##1\_marksp#1##2\_marksp##3\_nil{\CF_stripsp_ii##3##1##2\__nil#1\__nil\_nil}%
	\long\def\CF_stripsp_ii##1#1\__nil##2\_nil{\CF_stripsp_iii##1##2\_nil}%
	\long\def\CF_stripsp_iii##1##2\__nil##3\_nil##4{\unexpanded{##4{##2}}}%
	\long\def\CF_striplastsp##1##2{\expanded{\CF_striplastsp_i\_marksp##1\__nil#1\__nil\_nil{##2}}}%
	\long\def\CF_striplastsp_i##1#1\__nil##2\_nil{\CF_stripsp_iii##1\__nil\_nil}%
}\CF_stripsp{ }
\edef\CFhash{\string#}
\begingroup
	\catcode`\%12
	\gdef\CF_percent{%}
	\catcode`\_8
	\expandafter\gdef\csname CF\string_underscore\endcsname{_}
\endgroup
\def\CF_earg#1#2{\expandafter#1\expandafter{#2}}
\def\CF_etwoargs#1#2#3{\expandafter#1\expanded{{\unexpanded\expandafter{#2}}{\unexpanded\expandafter{#3}}}}
\def\CF_esecond_a#1#2{#2{#1}}
\def\CF_esecond#1#2{\expandafter\CF_esecond_a\expandafter{#2}{#1}}% \CF_esecond<{arg1>}{<arg2>} donne "<arg1>{*<arg2>}"
\def\CF_eesecond#1#2{\expandafter\expandafter\expandafter\CF_esecond_a\expandafter\expandafter\expandafter{#2}{#1}}% \CF_eesecond{<arg1>}{<arg2>} donne "<arg1>{**<arg2>}"
\def\CF_eafter_a#1#2{#2#1}
\def\CF_eafter#1#2{\expandafter\CF_eafter_a\expandafter{#2}{#1}}% \CF_eafter{<arg1>}{<arg2>} donne "<arg1>*<arg2>"
\def\CF_eeafter#1#2{\expandafter\expandafter\expandafter\CF_eafter_a\expandafter\expandafter\expandafter{#2}{#1}}% \CF_eeafter{<arg1>}{<arg2>} donne "<arg1>**<arg2>"
\def\CF_addtomacro#1#2{\expandafter\def\expandafter#1\expandafter{#1#2}}
\def\CF_eaddtomacro#1#2{\edef#1{\unexpanded\expandafter{#1}\unexpanded\expandafter{#2}}}
\def\CF_xaddtomacro#1#2{\edef#1{\unexpanded\expandafter{#1}#2}}
\def\CF_addtotoks#1#2{#1\expandafter{\the#1#2}}
\def\CF_eaddtotoks#1#2{\expandafter\CF_addtotoks\expandafter#1\expandafter{#2}}
\def\CF_assigntonil#1#2\_nil{\def#1{#2}}
\def\CF_ifnexttok#1#2#3{%
	\let\CF_ifnexttoktok=#1% <- espace indésirable, bugfix v1.31
	\def\CF_ifnexttokcodetrue{#2}%
	\def\CF_ifnexttokcodefalse{#3}%
	\futurelet\CF_temptok\CF_ifnexttoka
}
\def\CF_ifnexttoka{%
	\CF_ifx\CF_temptok\CF_sptoken
	{%
		\CF_ifnexttokb
	}
	{%
		\CF_ifx\CF_temptok\CF_ifnexttoktok
			\CF_ifnexttokcodetrue
			\CF_ifnexttokcodefalse
	}%
}
\expandafter\def\expandafter\CF_ifnexttokb\space{\futurelet\CF_temptok\CF_ifnexttoka}
\def\CF_ifstar#1{\CF_ifnexttok*{\CF_execfirst{#1}}}
\def\CF_testopt#1#2{\CF_ifnexttok[{#1}{#1[{#2}]}}
\def\CF_ifinteger#1{%
	\begingroup
		\afterassignment\CF_afterinteger
		\CF_cnt_cyclebonds0#1\relax
}
\def\CF_afterinteger#1\relax{%
	\endgroup
	\CF_ifempty{#1}%
}
\def\CF_iffirsttokis#1{% #1 commence-t-il par le token #2 (à lire) ?
	\futurelet\CF_toksa\CF_gobtonil#1\relax\_nil
	\CF_ifx\CF_toksa
}
\def\CF_iffirsttokin#1{% teste si le token qui commence #1 appartient aux tokens mis dans #2
	\futurelet\CF_toksa\CF_gobtonil#1\relax\_nil
	\CF_iffirsttokin_a
}
\def\CF_iffirsttokin_a#1{%
	\CF_ifempty{#1}
	{%
		\CF_execsecond
	}
	{%
		\futurelet\CF_toksb\CF_gobtonil#1\relax\_nil
		\CF_ifx\CF_toksa\CF_toksb
		{%
			\CF_execfirst
		}
		{%
			\CF_earg\CF_iffirsttokin_a{\CF_gobarg#1}%
		}%
	}%
}
\def\CF_ifinstr#1#2{%
	\def\CF_ifinstr_a##1#2##2\_nil{%
		\ifx\empty##2\empty
			\expandafter\CF_execsecond
		\else
			\expandafter\CF_execfirst
		\fi}%
	\CF_ifinstr_a#1\__nil#2\_nil
}
\def\CF_afterspace#1 #2\_nil{#2}
\def\CF_ifxcase_endif{\CF_ifxcase_endif}
\def\CF_ifxcase_elseif{\CF_ifxcase_elseif}
\def\CF_idto_endif#1\CF_ifxcase_endif{#1}
\def\CF_firstto_endif#1#2\CF_ifxcase_endif{#1}
\def\CF_ifxcase#1#2{% #1=token à tester   #2=token courant auquel #1 est comparé
	\CF_ifx\CF_ifxcase_elseif#2%
	{%
		\CF_idto_endif
	}
	{%
		\CF_ifx\CF_ifxcase_endif#2%
		{}
		{%
			\CF_ifx#1#2
				{\CF_firstto_endif}
				{\CF_ifxcase_a#1}%
		}%
	}%
}
\def\CF_ifxcase_a#1#2{\CF_ifxcase{#1}}

%--------------------------- Substitution ----------------------------
\def\CF_ifstartwith#1#2{% #1=<texte>  #2=<motif>
	\CF_ifempty{#1}%
	{%
		\CF_execsecond
	}
	{%
		\def\CF_startwithcode{#1}%
		\def\CF_startwithpattern{#2}%
		\CF_ifstartwitha
	}%
}
\def\CF_ifstartwitha{%
	\CF_grabfirstarg\CF_startwithcode\CF_firstargcode
	\CF_grabfirstarg\CF_startwithpattern\CF_firstargpattern
	\CF_ifx\CF_firstargcode\CF_firstargpattern
	{%
		\CF_earg\CF_ifempty\CF_startwithpattern
		{%
			\CF_execfirst
		}
		{%
			\CF_earg\CF_ifempty\CF_startwithcode
				\CF_execsecond
				\CF_ifstartwitha
		}%
	}
	{%
		\CF_execsecond
	}%
}
\def\CF_grabfirstarg#1#2{%
	\CF_ifx#1\empty
	{%
		\let#2\empty
	}
	{%
		\def\CF_grabmacro{#2}%
		\CF_earg\CF_ifbracefirst#1%
		{%
			\expandafter\CF_grabbracearg#1\_nil#1%
		}
		{%
			\CF_eafter{\futurelet\CF_nexttok\CF_grabfirstarga}#1\_nil#1%
		}%
	}%
}
\def\CF_grabfirstarga{%
	\CF_ifx\CF_nexttok\CF_sptoken
		\CF_grabspacearg
		\CF_grabnormalarg
}
\def\CF_grabbracearg#1{%
	\expandafter\def\CF_grabmacro{{#1}}%
	\CF_grabargassigntonil\relax
}
\expandafter\def\expandafter\CF_grabspacearg\space{%
	\expandafter\def\CF_grabmacro{ }%
	\CF_grabargassigntonil\relax
}
\def\CF_grabnormalarg#1{%
	\expandafter\def\CF_grabmacro{#1}%
	\CF_grabargassigntonil\relax
}
\def\CF_grabargassigntonil#1\_nil#2{\expandafter\def\expandafter#2\expandafter{\CF_gobarg#1}}
\def\CF_ifbracefirst#1{\CF_ifnum{\catcode\expandafter\expandafter\expandafter`\expandafter\CF_firsttonil\detokenize{#1.}\_nil=1 }}
\def\CF_substonly#1#2{% #1=entier maxi>0  #2=macro : dans la sc#1, remplace tous les <motif> par <pattern> sauf lorsque le motif est suivi d'un caractère >#1
	\def\CF_atendsubstitute{\edef#2{\the\CF_toks_subst}}% macro exécutée à la fin
	\let\CF_substnogroups\CF_substnogrouponly
	\CF_ifnum{#1>0 }
	{%
		\let\CF_testifx\empty
		\foreach\CF_x in {1,...,#1}
		{%
			\xdef\CF_testifx{\unexpanded\expandafter{\CF_testifx}\unexpanded\expandafter{\expandafter\ifx\CF_x\CF_nexttok1\fi}}%
		}%
		\let\CF_testif\empty
		\foreach\CF_x in {1,...,#1}
		{%
			\xdef\CF_testif{\unexpanded\expandafter{\CF_testif}\unexpanded\expandafter{\expandafter\if\CF_x\CF_nexttok1\fi}}%
		}%
		\CF_earg\CF_substi#2%
	}
	{%
		\CF_substall#2%
	}%
}
\def\CF_substnogrouponly{%
	\CF_etwoargs\CF_ifstartwith\CF_substcode\CF_substsubst
	{%
		\CF_grabfirstarg\CF_substcode\CF_temp
		\CF_eafter{\futurelet\CF_nexttok\CF_gobtonil}\CF_substcode\relax\_nil
		\CF_ifnum{0\CF_testifx=1 }% si le prochain token est \let-égal à 1...#1
		{%
			\edef\CF_nexttok{\expandafter\expandafter\expandafter\CF_firsttonil\expandafter\string\CF_substcode\_nil}% le détokéniser
			\CF_ifnum{0\CF_testif=1 }
			{%
				\CF_eaddtotoks\CF_toks_subst\CF_temp
				\CF_grabfirstarg\CF_substcode\CF_temp
				\CF_eaddtotoks\CF_toks_subst\CF_temp
			}
			{%
				\CF_eaddtotoks\CF_toks_subst\CF_substpattern
			}%
		}
		{%
			\CF_eaddtotoks\CF_toks_subst\CF_substpattern
		}%
		\CF_substgroups
	}
	{%
		\CF_earg\CF_ifempty\CF_substcode
		{%
			\CF_atendsubstitute
		}
		{%
			\CF_grabfirstarg\CF_substcode\CF_substauxarg
			\CF_eaddtotoks\CF_toks_subst\CF_substauxarg
			\CF_substgroups
		}%
	}%
}
\def\CF_substall#1{% #1=macro
	\def\CF_atendsubstitute{\edef#1{\the\CF_toks_subst}}% macro exécutée à la fin
	\let\CF_substnogroups\CF_substnogroupall
	\CF_earg\CF_substi#1%
}
\def\CF_substnogroupall{%
	\CF_etwoargs\CF_ifstartwith\CF_substcode\CF_substsubst
	{%
		\CF_eaddtotoks\CF_toks_subst\CF_substpattern
		\CF_grabfirstarg\CF_substcode\CF_temp
		\CF_substgroups
	}
	{%
		\CF_earg\CF_ifempty\CF_substcode
		{%
			\CF_atendsubstitute
		}
		{%
			\CF_grabfirstarg\CF_substcode\CF_substauxarg
			\CF_eaddtotoks\CF_toks_subst\CF_substauxarg
			\CF_substgroups
		}%
	}%
}
\def\CF_substi#1#2#3{% #1=<texte> #2=<motif> #3=<motif de substi>
	\def\CF_substcode{#1}\def\CF_substsubst{#2}\def\CF_substpattern{#3}%
	\CF_toks_subst={}%
	\CF_substgroups
}

\def\CF_substgroups{%
	\CF_earg\CF_ifbracefirst\CF_substcode
	{%
		\CF_grabfirstarg\CF_substcode\CF_substauxarg
		\begingroup
			\def\CF_atendsubstitute{%
				\expandafter\endgroup\expandafter\CF_addtotoks\expandafter\CF_toks_subst\expandafter{\expandafter{\the\CF_toks_subst}}%
				\CF_substgroups
			}%
			\CF_toks_subst{}% initialiser à vide
			\expandafter\def\expandafter\CF_substcode\CF_substauxarg
			\CF_substgroups
	}%
	{%
		\CF_substnogroups
	}%
}

%---------------------------- Paramètres -----------------------------
\def\CF_defifempty#1#2#3{\CF_ifempty{#2}{\def#1{#3}}{\def#1{#2}}}
\defKV[chemfig]{%
	atom style          = \def\CF_atomstyle                    {#1},
	baseline            = \CF_defifempty\CF_baseline           {#1}{0pt},% nouveau v 1.6e
	chemfig style       = \def\CF_chemfigstyle                 {#1},
	cram width          = \CF_defifempty\CF_crambasewidth      {#1}{1.5ex},
	cram dash width     = \CF_defifempty\CF_cramdashlength     {#1}{1pt},
	cram dash sep       = \CF_defifempty\CF_cramdashsep        {#1}{2pt},
	atom sep            = \CF_defifempty\CF_atomsep            {#1}{3em},
	bond offset         = \CF_defifempty\CF_bondoffset         {#1}{2pt},
	double bond sep     = \CF_defifempty\CF_doublesep          {#1}{2pt},
	angle increment     = \CF_defifempty\CF_angleincrement     {#1}{45},
	node style          = \def\CF_nodestyle                    {#1},
	bond style          = \def\CF_bondstyle                    {#1},
	cycle radius coeff  = \CF_defifempty\CF_cycleradiuscoeff   {#1}{0.75},
	stack sep           = \CF_defifempty\CF_stacksep           {#1}{1.5pt},
	compound style      = \def\CF_defaultcompoundstyle         {#1},
	compound sep        = \CF_defifempty\CF_compoundsep        {#1}{5em},
	arrow offset        = \CF_defifempty\CF_arrowoffset        {#1}{1em},
	arrow angle         = \CF_defifempty\CF_arrowangle         {#1}{0},
	arrow coeff         = \CF_defifempty\CF_arrowlength        {#1}{1},
	arrow style         = \def\CF_defaultarrowstyle            {#1},
	arrow double sep    = \CF_defifempty\CF_arrowdoublesep     {#1}{2pt},
	arrow double coeff  = \CF_defifempty\CF_arrowdoubleposstart{#1}{0.6},
	arrow label sep     = \CF_defifempty\CF_arrowlabelsep      {#1}{3pt},
	arrow head          = \CF_defifempty\CF_arrowhead          {#1}{-CF},
	init anchor         = \CF_defifempty\CF_nextnodeanchor     {#1}{text},% new v1.71
	+ sep left          = \CF_defifempty\CF_signspaceante      {#1}{0.5em},
	+ sep right         = \CF_defifempty\CF_signspacepost      {#1}{0.5em},
	+ vshift            = \CF_defifempty\CF_signvshift         {#1}{0pt},
	harrow minwidth     = \CF_defifempty\CF_harrow_min_width   {#1}{3em},
	label xsep          = \CF_defifempty\CF_harrow_label_xsep  {#1}{3pt},
	label align         = \CF_set_label_align                  {#1},
	name sep            = \CF_defifempty\CF_hreac_name_sep     {#1}{1.5ex},
	hreac anchor        = \def\CF_hreac_init_anchor            {#1},
	hreac sep           = \CF_defifempty\CF_hreac_sep          {#1}{0.5em},
	save                = \testboolKV{#1}
                            {\def\CF_save{\immediate\write\CF_write}\relax}
                            {\let\CF_save\CF_gobarg},
	save name           = \set_save_name\CF_save_name{#1},
}
\def\set_save_name#1#2{% #1=macro recevant le nom  #2=argument
	\CF_ifinstr{#2}*
		{\set_save_name_a#2\_nil#1}
		{\def#1{#2}}%
}
\def\set_save_name_a#1*#2\_nil#3{%
	\def#3{#1\number\CF_cnt_save#2}%
}
\def\CF_set_label_align#1{%
	\let\CF_harrow_label_align_right\hss
	\let\CF_harrow_label_align_left \hss
	\expandafter\CF_ifxcase\CF_firsttonil#1c\_nil
		c {}
		l {\let\CF_harrow_label_align_left\empty}
		r {\let\CF_harrow_label_align_right\empty}
	\CF_ifxcase_elseif
		\CF_error{Illegal alignment directive '#1'}%
	\CF_ifxcase_endif
}
\def\setchemfig{\setKV[chemfig]}
\def\resetchemfig{\restoreKV[chemfig]}
\setKVdefault[chemfig]{%
	atom style          = {},% code tikz mis à la fin de every node/.style
	use atom strut      = false,% met à chaque atome le strut de l'atome précédent nouveau 1.7
	baseline            = {0pt},% nouveau v 1.6e
	chemfig style       = {},% code tikz mis à la fin de l'arugment optionnel de \tikzpicture
	bond join           = false,
	fixed length        = false,
	cram rectangle      = false,
	cram width          = 1.5ex,
	cram dash width     = 1pt,
	cram dash sep       = 2pt,
	atom sep            = 3em,
	bond offset         = 2pt,
	double bond sep     = 2pt,
	angle increment     = 45,
	node style          = {},
	bond style          = {},
	cycle radius coeff  = 0.75,
	stack sep           = 1.5pt,
	autoreset cntcycle  = true,
	show cntcycle       = false,
	debug               = false,
	scheme debug        = false,
	compound style      = {},
	compound sep        = 5em,
	arrow offset        = 1em,
	arrow angle         = 0,
	arrow coeff         = 1,
	arrow style         = {},
	arrow double sep    = 2pt,
	arrow double coeff  = 0.6,
	arrow double harpoon= true,
	arrow label sep     = 3pt,
	arrow head          = -CF,
	init anchor         = text,  % new v1.71
	+ sep left          = 0.5em,
	+ sep right         = 0.5em,
	+ vshift            = 0pt,
	gchemname           = true,  % ajout de cette clé v1.6c
	schemestart code    = {},    % ajout de cette clé v1.6c
	schemestop code     = {},    % ajout de cette clé v1.6c
	harrow minwidth     = 3em,   % dimension minimale de la flèche
	hreac debug         = false,
	label xsep          = 3pt,   % dist H supplémentaire
	label align         = c,     % vaut c, l ou r
	name sep            = 1.5ex,   % distance v entre composé et nom
	hreac anchor        = text,  % ancre de la ligne de base du 1er composé
	hreac sep           = 0.5em, % espace horizontal entre composés
	save                = false,
	save name           = \jobname-fig*,
}%

%%%%%%%%%%%%%%%%%%%%%%%%%%%%%%%%%%%%%%%%%%%%%%%%%%%%%%%%%%%%%%%%%%%%%%
%%%               D E S S I N     M O L É C U L E S                %%%
%%%%%%%%%%%%%%%%%%%%%%%%%%%%%%%%%%%%%%%%%%%%%%%%%%%%%%%%%%%%%%%%%%%%%%
\def\CF_sanitizecatcode{%
	\CF_setcatcodes{\[12,\]12,\:12,\(12,\)12,\,12,\-12,\=12,\~12,\!12,\?12,\<12,\>12,\;12,\*12,\|12,\#12,\@12}%
}

\def\printatom#1{\ifmmode\rm#1\else$\rm#1$\fi}

\def\chemskipalign{%
	\CF_doifempty\CF_bondoutcontentsaved% sauf si un \chemskipalign a été fait à l'atome précédent
	{%
		\global\let\CF_bondoutcontentsaved\CF_bondoutcontent% sauvegarder l'atome d'où vient la liaison
	}%
	\let\CF_nodestrut\empty
}

\def\definesubmol{\CF_define_submoltrue\CF_defsubmol}
\def\redefinesubmol{\CF_define_submolfalse\CF_defsubmol}

\def\CF_defsubmol#1{%
	\CF_cnt_atomgroup=0 % nombre d'arguments supposé
	\def\CF_temp{#1}% nom
	\futurelet\CF_toksa\CF_submoltestnxttok
}

\def\CF_submoltestnxttok{%
	\if[\noexpand\CF_toksa\CF_test_ok\fi\CF_execsecond
	{%
		\begingroup\CF_sanitizecatcode\CF_earg\CF_submolgrabopt{\CF_temp}%
	}
	{%
		\afterassignment\CF_submoltestnxttok_a% pas d'argument entre crochet
		\CF_cnt_atomgroup=0% cherche le nombre d'arguments éventuels
	}%
}

\def\CF_submoltestnxttok_a{\futurelet\CF_toksa\CF_submoltestnxttok_b}

\def\CF_submoltestnxttok_b{%
	\if[\noexpand\CF_toksa\CF_test_ok\fi\CF_execsecond
	{%
		\begingroup\CF_sanitizecatcode\CF_earg\CF_submolgrabopt{\CF_temp}%
	}
	{%
		\CF_earg\CF_defsubmol_a\CF_temp{}%
	}%
}

\def\CF_submolgrabopt#1[#2]#{\endgroup\CF_defsubmol_a{#1}{#2}}

\def\CF_defsubmol_a#1{% #1 nom
	\CF_ifnum{0\CF_ifnum{\CF_cnt_atomgroup<0 }1{\CF_ifnum{\CF_cnt_atomgroup>9 }10}>0 }
	{%
		\CF_error{Invalid number of arguments in submol \detokenize\expandafter{\string#1}. Defining it with 0 argument}%
		\CF_cnt_atomgroup=0
	}
	{}%
	\ifcat\relax\expandafter\noexpand\CF_firsttonil#1\_nil\CF_test_ok\fi\CF_execsecond% si #1 est une séquence de contrôle
	{%
		\expandafter\ifdefined\CF_firsttonil#1\_nil
			\ifCF_define_submol
				\CF_warning{the submol \expandafter\string\CF_firsttonil#1\_nil\space is already defined, the previous definition is lost}%
			\fi
		\fi
		\begingroup
			\CF_sanitizecatcode
			\CF_defsubmol_b{#1}%
	}
	{%
		\ifcsname CF__#1\endcsname
			\ifCF_define_submol
				\CF_warning{the submol "#1" is already defined, the previous definition is lost}%
			\fi
		\fi
		\begingroup
			\CF_sanitizecatcode
			\expandafter\CF_defsubmol_b\csname CF__#1\endcsname
	}%
}

\def\CF_defsubmol_b#1#2#3{% #1 nom sous forme de macro, #2 = code si liaison arrive de droite , #3 = code si liaison arrive de gauche, \CF_cnt_atomgroup = nombre d'arguments
	\def\CF_tempa{#2}%
	\CF_doifnotempty{#2}
	{%
		\CF_etwoargs\CF_ifinstr{\detokenize{#2}}\CFhash
		{%
		\CF_esecond{\CF_earg\CF_substonly{\number\CF_cnt_atomgroup}\CF_tempa}\CFhash\CFhash
		}
		{}%
	}%
	\def\CF_tempb{#3}%
	\CF_etwoargs\CF_ifinstr{\detokenize{#3}}\CFhash
	{%
	\CF_esecond{\CF_earg\CF_substonly{\number\CF_cnt_atomgroup}\CF_tempb}\CFhash\CFhash
	}
	{}%
	\CF_esecond{\CF_esecond{\CF_defsubmolc{#1}}\CF_tempa}\CF_tempb
}

\def\CF_defsubmolc#1#2#3{% #1 nom sous forme de macro, #2 = code si liaison arrive de droite , #3 = code si liaison arrive de gauche, \CF_cnt_atomgroup = nombre d'arguments
	\endgroup
	\begingroup
		\toks0{\expandafter\gdef\csname CF_thesubmol\endcsname}%
		\CF_defsubmol_buildpreamble1\CF_cnt_atomgroup
		\CF_sanitizecatcode
		\catcode`\#6
		\endlinechar-1
		\everyeof{\noexpand}%
		\CF_ifempty{#2}%
		{%
			\scantokens\expandafter{\the\toks0{#3}}%
		}%
		{%
			\scantokens\expandafter
			{%
				\the\toks0
				{%
					\romannumeral% bugfix 1.52
						\csname CF_exec%
							\ifdim\csname CF_currentangle\endcsname pt>90pt
								\ifdim\csname CF_currentangle\endcsname pt<270pt
									first%
								\else
									second%
								\fi
							\else
								second%
							\fi
						\endcsname
							{0 #2}{0 #3}%
				}%
			}%
		}%
	\endgroup
	\let#1\CF_thesubmol
}

\def\CF_defsubmol_buildpreamble#1#2{% #1=itération courante  #2=nb max d'itérations
	\CF_ifnum{#1>#2}
	{}
	{%
		\toks0\expandafter{\the\toks0\CFhash#1}%
		\CF_earg\CF_defsubmol_buildpreamble{\the\numexpr#1+1}{#2}%
	}%
}

\def\CF_searchnode#1#2#3{% cherche un noeud au début de #1 l'assigne dans la sc #2 et met le reste dans #3
	\let#2\empty
	\def#3{#1}%
	\CF_iffirsttokis{#1}\CF_sptoken
	{%
		\CF_earg\CF_searchnode_a{\romannumeral-`\.\noexpand#1}#2#3% ignore les espaces au début du groupe d'atome
	}%
	{%
		\CF_searchnode_a{#1}#2#3%
	}%
}

\def\CF_searchnode_a#1#2#3{%
	\CF_ifempty{#1}%
	{%
		\let#3\empty
	}
	{%
		\CF_iffirsttokis{#1}\CF_sptoken
		{%
			\CF_addtomacro#2{ }%
			\CF_earg\CF_searchnode_a{\CF_afterspace#1\_nil}#2#3%
		}%
		{%
			\CF_ifx\CF_toksa\bgroup
			{%
				\CF_eaddtomacro#2{\expandafter{\CF_firsttonil#1\_nil}}%
				\CF_earg\CF_searchnode_a{\CF_gobarg#1}#2#3%
			}%
			{%
				\CF_ifx!\CF_toksa% Bugfix v1.5
				{%
					\def\CF_seeksubmloltemp{#1}%
					\CF_searchsubmola
					\CF_earg\CF_searchnode_a\CF_seeksubmloltemp#2#3%
				}
				{%
					\CF_iffirsttokin_a{-=(*<>~}%
					{%
						\def#3{#1}%
					}%
					{%
						\CF_eaddtomacro#2{\CF_firsttonil#1\_nil}%
						\CF_earg\CF_searchnode_a{\CF_gobarg#1}#2#3%
					}%
				}
			}%
		}%
	}%
}

% on sait que #1 commence par -,=,~,<,>. On analyse cette liaison
% #2 reçoit le type de liaisons (1 pour -, 2 pour =, 3 pour ~)
\def\CF_assignbondcode#1#2{%
	\futurelet\CF_toksa\CF_gobtonil#1\_nil
	\edef#2{%
		\ifx-\CF_toksa1\else
		\ifx=\CF_toksa2\else
		\ifx~\CF_toksa3\else
		\ifx>\CF_toksa4\else
		\ifx<\CF_toksa5\else0% si 0 --> il y a une erreur non due à l'utilisateur
		\fi\fi\fi\fi\fi
		}%
	\ifnum#2>3 % si c'est une liaison de Cram
		\CF_earg\CF_iffirsttokis{\CF_gobarg#1\relax}:% teste le token suivant
		{%
			\edef#2{\number\numexpr#2+2}% si c 'est un ":", signe du pointillé, ajoute 2
		}%
		{%
			\CF_ifx|\CF_toksa% si c 'est un "|", signe du triangle évidé, ajouter 4
			{%
				\edef#2{\number\numexpr#2+4}%
			}
			{}%
		}%
	\fi
}

\def\CF_grabbondoffset_a#1,#2\_nil{%
	\def\CF_startoffset{#1}\def\CF_endoffset{#2}%
}

\def\CF_grabbondoffset#1(#2)#3\_nil{%
	\CF_doifnotempty{#2}%
	{%
		\CF_ifinstr{#2},%
		{%
			\CF_grabbondoffset_a#2\_nil
		}%
		{%
			\def\CF_startoffset{#2}%
		}%
	}%
	\def\CF_remainafterbond{#3}%
}

\def\CF_analysebond#1#2{%
	\CF_assignbondcode{#1}#2%
	\expandafter\def\expandafter\CF_remainafterbond\expandafter{\CF_gobarg#1}%mange le premier signe de la liaison
	\let\CF_doublebondtype\CF_zero
	\ifnum#2=2 % si c'est une double liaison, regarde s'il y a un + ou - derrière
		\CF_eafter{\futurelet\CF_toksa\CF_gobtonil}{\CF_gobarg#1\_nil}%
		\ifx^\CF_toksa
			\def\CF_doublebondtype{1}%
			\expandafter\def\expandafter\CF_remainafterbond\expandafter{\CF_gobtwoargs#1}% mange le "^"
		\else
			\expandafter\ifx\CF_underscore\CF_toksa
				\def\CF_doublebondtype{2}%
				\expandafter\def\expandafter\CF_remainafterbond\expandafter{\CF_gobtwoargs#1}% mange le "_"
			\fi
		\fi
	\else
		\ifnum#2>5 % si c'est une laision de Cram pointillée ou triangle évidé
			\expandafter\def\expandafter\CF_remainafterbond\expandafter{\CF_gobtwoargs#1}% mange un caractère de plus
		\fi
	\fi
	\CF_earg\CF_iffirsttokis\CF_remainafterbond\CFhash
	{%
		\CF_eesecond\CF_iffirsttokis{\expandafter\CF_gobarg\CF_remainafterbond.}(%si parenthèse juste après
		{%
			\expandafter\CF_grabbondoffset\CF_remainafterbond\_nil
		}
		{}%
	}%
	{}%
	\CF_earg\CF_iffirsttokis\CF_remainafterbond @%
	{%
		\expandafter\CF_grabmovearg\CF_remainafterbond\_nil
	}%
	{}%
	\CF_earg\CF_iffirsttokis{\CF_remainafterbond}[%
	{%
		\expandafter\CF_analyseoptarg\CF_remainafterbond\_nil\CF_remainafterbond
	}%
	{%
		\let\CF_currentstringangle\CF_defaultstringangle
		\let\CF_currentlength\CF_defaultlength
		\let\CF_currentfromatom\CF_defaultfromatom
		\let\CF_currenttoatom\CF_defaulttoatom
		\let\CF_currenttikz\CF_defaulttikz
		\let\CF_movebondname\empty
	}%
	\ifCF_in_cycle
		\pgfmathsetmacro\CF_cycleincrementangle{360/\CF_cyclenum+\CF_initcycleangle}%
		\edef\CF_currentstringangle{::+\CF_cycleincrementangle}%
		\def\CF_initcycleangle{0}%
		\let\CF_currentlength\CF_defaultlength% et on ignore la longueur de liaison spécifiée
	\fi
	\CF_earg\CF_setbondangle{\CF_currentstringangle}\CF_currentangle
}

\def\CF_setbondangle#1#2{% le code de la direction est contenu dans #1, en sortie, #2 contient l'angle
	\CF_ifempty{#1}%
	{%
		\let#2\CF_defaultangle
	}
	{%
		\if:\expandafter\noexpand\CF_firsttonil#1\_nil
			\if:\expandafter\expandafter\expandafter\noexpand\expandafter\CF_firsttonil\CF_gobarg#1\_nil
				\pgfmathsetmacro#2{\CF_previousangle+\expandafter\CF_gobarg\CF_gobarg#1}%
			\else
				\pgfmathsetmacro#2{\CF_gobarg#1}%
			\fi
		\else
			\pgfmathsetmacro#2{#1*\CF_angleincrement}%
		\fi% puis normalise l'angle entre 0 et 360
		\ifdim\ifdim#2pt<0pt -\fi#2pt>360pt
			\pgfmathsetmacro#2{#2-360*floor(#2/360)}%
		\fi% si |#2|>360
		\ifdim#2pt<0pt
			\pgfmathsetmacro#2{#2+360}%
		\fi
	}%
}

\def\CF_analysemovearg#1,#2\_nil#3{%
	\def#3{#1}%
	\def\CF_movebondcoeff{#2}%
}

% Argument limités légitimes ici car #2 (qui est ce qui suit "@{<nom>}" dans l'argument optionnel) ne DOIT PAS commencer par une accolade.
\def\CF_grabmovearg @#1#2\_nil{%
	\CF_ifinstr{#1},%
	{%
		\CF_analysemovearg#1\_nil\CF_movebondname
	}%
	{%
		\def\CF_movebondname{#1}%
		\def\CF_movebondcoeff{0.5}%
	}%
	\def\CF_remainoptarg{#2}%
}

\def\CF_testemptyandassign#1#2#3{%
	\CF_ifempty{#2}
	{%
		\let#1#3%
	}
	{%
		\def#1{#2}%
	}%
}

\def\CF_parseoptlist#1,#2,#3,#4,#5\_nil{%
	\CF_testemptyandassign\CF_currentstringangle{#1}\CF_defaultstringangle
	\CF_testemptyandassign\CF_currentlength     {#2}\CF_defaultlength
	\CF_testemptyandassign\CF_currentfromatom   {#3}\CF_defaultfromatom
	\CF_testemptyandassign\CF_currenttoatom     {#4}\CF_defaulttoatom
	\expandafter\CF_testemptyandassign\expandafter\CF_currenttikz\expandafter{\CF_sanitizelastitem#5,\empty\_nil}\CF_defaulttikz
}%

\def\CF_analyseoptarg[#1]{%
	\CF_doifnotempty{#1}%
	{%
		\CF_iffirsttokis{#1}@%
		{%
			\CF_grabmovearg#1\_nil
		}
		{%
			\let\CF_movebondname\empty
			\def\CF_remainoptarg{#1}%
		}%
		\expandafter\CF_parseoptlist\CF_remainoptarg,\empty,\empty,\empty,\empty\_nil
	}%
	\CF_analyseoptarg_a\relax
}

\def\CF_analyseoptarg_a#1\_nil#2{%
	\expandafter\def\expandafter#2\expandafter{\CF_gobarg#1}%
}

\def\CF_searchsubmol#1#2{% cherche et remplace ! au début de #1. #1=code #2=macro recevant le résultat
	\def\CF_seeksubmloltemp{#1}%
	\CF_searchsubmola
	\let#2\CF_seeksubmloltemp
}

\def\CF_searchsubmola{%
	\expandafter\def\expandafter\CF_seeksubmloltemp\expandafter{\romannumeral-`\.\expandafter\noexpand\CF_seeksubmloltemp}%
	\CF_earg\CF_iffirsttokis\CF_seeksubmloltemp!%
	{%
		\CF_eesecond{\def\CF_seeksubmloltemp}{\expandafter\CF_gobarg\CF_seeksubmloltemp}% enlève le "!"
		\CF_ifx\empty\CF_seeksubmloltemp
		{%
			\CF_error{no submol name found after "!"}%
		}
		{}%
		\ifcat\relax\expandafter\expandafter\expandafter\noexpand\expandafter\CF_firsttonil\CF_seeksubmloltemp*\_nil
			\expandafter\CF_searchsubmol_b\CF_seeksubmloltemp\_nil
		\else
			\expandafter\CF_searchsubmol_c\CF_seeksubmloltemp\_nil
		\fi
		\CF_searchsubmola
	}%
	{}%
}

\def\CF_searchsubmol_b#1{\CF_searchsubmol_d#1\relax}

\def\CF_searchsubmol_c#1{\expandafter\CF_searchsubmol_d\csname CF__#1\endcsname\relax}% nom de la sous molécule

\def\CF_searchsubmol_d#1#2\_nil{% #1=macro de la sous molécule #2=reste du code commençant par \relax
	\edef\CF_seeksubmloltemp{\unexpanded\expandafter\expandafter\expandafter\expandafter\expandafter\expandafter\expandafter{\expandafter#1\CF_gobarg#2}}% supprime le \relax puis ajoute la macro au début et la 2-développe
}

\def\CF_insertemptygroup#1{% insère {} au début de la sc #1
	\edef#1{{}\unexpanded\expandafter{#1}}%
}
\def\chemfig{\CF_testopt\CF_chemfig_a{}}

\def\CF_chemfig_a[#1]{%
	\begingroup
		\CF_sanitizecatcode
		\CF_etwoargs\CF_chemfig_b\CF_begintikzpicture\CF_endtikzpicture[#1]%
}

\def\CF_chemfig_b#1#2[#3]#4{%
	\endgroup
	\begingroup
		\setchemfig{#3}%
		\ifboolKV[chemfig]{save}
		{%
			\global\advance\CF_cnt_save1
			\immediate\openout\CF_write=\CF_save_name.tex\relax
		}
		{}%
		\CF_ifinsidetikz
		{%
			\pgfinterruptpicture
			\CF_save{\noexpand\pgfinterruptpicture}%
			\def\CF_atendofchemfig{%
				\endpgfinterruptpicture
				\CF_save{\noexpand\endpgfinterruptpicture}%
			}%
		}
		{%
			\let\CF_atendofchemfig\empty
		}%
		\CF_save
		{%
		\unexpanded{#1}[%
						remember picture,%
						every node/.style={%
							anchor=base,%
							inner sep=0pt,%
							outer sep=0pt,%
							minimum size=0pt,%
							\unexpanded\expandafter{\CF_atomstyle}%
							},%
						baseline=\CF_baseline,%
						\unexpanded\expandafter{\CF_chemfigstyle}%
						]%
		}%
		\expanded
		{% début du tikzpicture
			\unexpanded{#1}[%
				remember picture,%
				every node/.style={%
					anchor=base,%
					inner sep=0pt,%
					outer sep=0pt,%
					minimum size=0pt,%
					\unexpanded\expandafter{\CF_atomstyle}%
					},%
				baseline=\CF_baseline,%
				\unexpanded\expandafter{\CF_chemfigstyle}%
				]%
		}%
			\begingroup% \endgroup rajouté en sortie de tracé par \CF_chemfig_d
				\let\CF_hooklist\empty
				\ifboolKV[chemfig]{fixed length}
				{%
					\CF_macro_fixed_bond_lengthtrue
				}
				{%
					\CF_macro_fixed_bond_lengthfalse
				}%
				\ifboolKV[chemfig]{bond join}
				{%
					\let\CF_drawaxisbond\CF_drawaxisbondjoin
				}
				{%
					\let\CF_drawaxisbond\CF_drawaxisbondnojoin
				}%
				\ifboolKV[chemfig]{cram rectangle}
				{%
					\let\CF_clipcramornot\CF_gobtikzinstruction
				}
				{%
					\let\CF_clipcramornot\clip
				}%
				\CF_in_cyclefalse
				\CF_cnt_group0
				\ifboolKV[chemfig]{autoreset cntcycle}
				{%
					\global\CF_cnt_cycle0
				}
				{}%
				\let\CF_lastaction\CF_zero% 0=début du dessin 1=tracé d'un noeud 2=tracé d'une liaison
				\let\CF_startoffset\empty
				\let\CF_endoffset\empty
				\let\CF_bondoutcontentsaved\empty
				\def\CF_cycleanglecorrection{180/\CF_cyclenum}%
				\def\CF_defaultangle{0}%
				\def\CF_defaultstringangle{:0}% angle pris par défaut si le champ est vide
				\def\CF_defaultlength{1}%
				\let\CF_defaultfromatom\empty% numero de l'atome d'où partent les liaisons par défaut
				\let\CF_defaulttoatom\empty% numéro de l'atome où arrivent les laisons par défaut
				\let\CF_defaulttikz\empty
				\let\CF_previousbondangle\empty
				\def\CF_previousbondtype{0}%
				\def\CF_previousnode{}%
				\def\CF_previous_offset{0pt}%
				\let\CF_joinbond\CF_zero
				\let\CF_previoustikz\empty
				\everyeof{\_nil}\endlinechar-1
				\CF_sanitizecatcode
				\expandafter\CF_assigntonil\expandafter\CF_molecule\scantokens{#4}%
				\expandafter\CF_substall\expandafter\CF_molecule\expandafter{\CFhash}\CFhash
				\CF_earg\CF_chemfig_c{\CF_molecule}%
			%\endgroup <-- rajouté par \CF_chemfig_d
			\CF_save{\unexpanded{#2}\CF_percent}%
		#2% fin du tikzpicture
		\CF_atendofchemfig
		\ifboolKV[chemfig]{save}
		{%
			\immediate\closeout\CF_write
		}
		{}%
	\endgroup
	\let\CF_flipstate\CF_zero
}

\def\CF_chemfig_c#1{% #1 est le code de la molécule
	\ifnum\CF_lastaction=3
		\ifCF_in_cycle
			\def\CF_defaultangle{0}%
		\else
			\ifnum\CF_cnt_cyclebonds=0 % si c'est le début d'un cycle
				\pgfmathsetmacro\CF_defaultangle{\CF_previousangle+180+\CF_cycleanglecorrection}% on met la liaison à +180° + correction
			\else
				\pgfmathsetmacro\CF_defaultangle{\CF_previousangle-90+180/\CF_cyclenum}% sinon à la bissectrice du sommet du cycle
			\fi
		\fi
		\let\CF_defaultstringangle\empty
	\fi
	\let\CF_currentangle\CF_defaultangle
	\def\CF_molecule{#1}%
	\CF_earg\CF_searchsubmol\CF_molecule\CF_molecule% alias en premier ?
	\if[\expandafter\expandafter\expandafter\noexpand\expandafter\CF_firsttonil\CF_molecule\_nil
		\expandafter\CF_analyseoptarg\CF_molecule\_nil\CF_molecule
		\CF_earg\CF_setbondangle{\CF_currentstringangle}\CF_currentangle
		\let\CF_defaultangle\CF_currentangle
		\let\CF_previousangle\CF_currentangle
		\CF_doifnotempty\CF_currentstringangle{\let\CF_defaultangle\CF_currentangle}%
		\CF_doifnotempty\CF_currentlength     {\let\CF_defaultlength\CF_currentlength}%
		\CF_doifnotempty\CF_currentfromatom   {\let\CF_defaultfromatom\CF_currentfromatom}%
		\CF_doifnotempty\CF_currenttoatom     {\let\CF_defaulttoatom\CF_currenttoatom}%
		\CF_doifnotempty\CF_currenttikz       {\let\CF_defaulttikz\CF_currenttikz}%
		\CF_earg\CF_searchsubmol\CF_molecule\CF_molecule
	\fi
	\edef\CF_defaultstringangle{:\CF_defaultangle}%
	\let\CF_currentlength\CF_defaultlength
	\let\CF_currentfromatom\CF_defaultfromatom
	\let\CF_currenttoatom\CF_defaulttoatom
	\let\CF_currenttikz\CF_defaulttikz
	\ifCF_in_cycle% si on commence un cycle
		\let\CF_currentangle\CF_previousangle
		\pgfmathsetmacro\CF_cycle_arcinitangle{\CF_currentangle+\CF_initcycleangle+180/\CF_cyclenum+90}%
		\pgfmathsetmacro\CF_centeroffset{\CF_currentlength*\CF_atomsep/(2*sin(180/\CF_cyclenum))}%
		\CF_save{\noexpand\node[at=(\CF_bondoutnode),shift=(\CF_cycle_arcinitangle:\CF_centeroffset pt),anchor=center](cyclecenter\number\CF_cnt_cycle){};}%
		\node[at=(\CF_bondoutnode),shift=(\CF_cycle_arcinitangle:\CF_centeroffset pt),anchor=center](cyclecenter\number\CF_cnt_cycle){};% le centre du cycle
		\ifboolKV[chemfig]{show cntcycle}
		{%
			\CF_save{\noexpand\node[at=(cyclecenter\number\CF_cnt_cycle),anchor=center,overlay]{\noexpand\tiny\number\CF_cnt_cycle};}%
			\node[at=(cyclecenter\number\CF_cnt_cycle),anchor=center,overlay]{\tiny\number\CF_cnt_cycle};%
		}
		{}%
		\ifCF_cycle_arc% on doit tracer l'arc de cercle dans le cycle ?
			\pgfmathsetmacro\CF_cycle_arcradius{\CF_cycleradiuscoeff*\CF_currentlength*\CF_atomsep/(2*tan(180/\CF_cyclenum))}%
			\CF_save{\noexpand\node[at=(cyclecenter\number\CF_cnt_cycle),shift=(\CF_cycle_arcstartangle:\CF_cycle_arcradius pt)](initarc){};}%
			\node[at=(cyclecenter\number\CF_cnt_cycle),shift=(\CF_cycle_arcstartangle:\CF_cycle_arcradius pt)](initarc){};% le début de l'arc
			\CF_save{\CF_eafter{\noexpand\draw[}\CF_cycle_arcdirecttikz](initarc) arc (\CF_cycle_arcstartangle:\CF_cycle_arcendangle:\CF_cycle_arcradius pt);}%
			\CF_eafter{\draw[}\CF_cycle_arcdirecttikz](initarc) arc (\CF_cycle_arcstartangle:\CF_cycle_arcendangle:\CF_cycle_arcradius pt);%
		\fi
	\else
		\let\CF_currentangle\CF_defaultangle
	\fi
	\ifnum\CF_lastaction=0
		\let\CF_previousangle\CF_defaultangle
		\CF_save{\noexpand\node(CF_node){};}%
		\node(CF_node){};
		\CF_earg\CF_iffirsttokin{\CF_molecule}{-=(*~?<>}%
		{%
			\CF_insertemptygroup\CF_molecule
		}%
		{}%
	\fi
	\CF_chemfig_d
}

\def\CF_chemfig_d{%
	\let\CF_nextaction\CF_chemfig_d% à priori, on reboucle
	\CF_ifx\CF_molecule\empty
	{%
		\let\CF_nextaction\endgroup
	}
	{%
		\CF_earg\CF_searchnode{\CF_molecule}\CF_currentatomgroup\CF_molecule
		\CF_ifx\empty\CF_currentatomgroup% pas de noeud pour commencer ?
		{%
			\def\CF_bondoutnode{%
				n\CF_lastgroupnumber-%
				\ifx\CF_currentfromatom\empty
					\ifdim\CF_currentangle pt<90pt
						\number\CF_cnt_atomgroup
					\else
						\ifdim\CF_currentangle pt>270pt
							\number\CF_cnt_atomgroup
						\else
							1%
						\fi
					\fi
				\else
					\CF_currentfromatom
				\fi
				}%
			\CF_eafter{\futurelet\CF_toksa\CF_gobtonil}{\CF_molecule\relax\_nil}%
			\CF_iffirsttokin_a{-=<>~}% la suite est une liaison
			{%
				\ifnum\CF_lastaction=2 % c'est la deuxième liaison consécutive ?
					\CF_insertemptygroup\CF_molecule% insère un groupe vide
					\edef\CF_bondoutnode{\CF_bondoutnode}%
				\else
					\ifCF_in_cycle
						\advance\CF_cnt_cyclebonds1
					\fi
					\CF_earg\CF_analysebond{\CF_molecule}\CF_bondtype
					\edef\CF_bondoutnode{\CF_bondoutnode}%
					\let\CF_molecule\CF_remainafterbond
					\ifCF_in_cycle
						\ifnum\CF_cnt_cyclebonds=\CF_cyclenum\relax
							\expandafter\expandafter\expandafter\CF_execfirst
						\else
							\ifnum\CF_cnt_cyclebonds=1
								\let\CF_cyclefirsttikz\CF_currenttikz
								\CF_doifnotempty\CF_startoffset{\let\CF_cyclejoinlast\CF_zero}%
							\fi
							\expandafter\expandafter\expandafter\CF_execsecond
						\fi
					\else
						\expandafter\CF_execsecond
					\fi
					{%
						\let\CF_nextaction\endgroup
						\CF_drawbond\CF_bondtype{\CF_bondoutnode}{\CF_hookcycle}\CF_previousatomgroup\CF_hookatomgroup
					}%
					{%
						\CF_save{\noexpand\node[at=(\CF_bondoutnode\ifCF_in_cycle\else\ifCF_macro_fixed_bond_length.\CF_currentangle\fi\fi),shift=(\ifcase\CF_flipstate\or180-\or-\fi\CF_currentangle:\CF_currentlength*\CF_atomsep)](CF_node){};}%
						\node[at=(\CF_bondoutnode\ifCF_in_cycle\else\ifCF_macro_fixed_bond_length.\CF_currentangle\fi\fi),shift=(\ifcase\CF_flipstate\or180-\or-\fi\CF_currentangle:\CF_currentlength*\CF_atomsep)](CF_node){};
						\let\CF_previousangle\CF_currentangle
						\def\CF_lastaction{2}% on a tracé une liaison
					}%
				\fi
				\ifcat\relax\detokenize\expandafter{\romannumeral-`\.\expandafter\noexpand\CF_molecule}\relax
				% s'il ne reste plus rien après la liaison (sans tenir compte de l'espace devant)-> insère un groupe vide
					\CF_insertemptygroup\CF_molecule
				\fi
			}%
			{%
				\edef\CF_bondoutnode{\CF_bondoutnode}% évalue le l'atome de départ de liaison
				\CF_ifx(\CF_toksa% une parenthèse pour commencer ?
				{%
					\ifnum\CF_lastaction=2 % il y avait une liaison juste avant ?
						\CF_insertemptygroup\CF_molecule
					\else
						\CF_earg\CF_grabsubmol{\CF_molecule}%
						\begingroup
							\ifCF_in_cycle\def\CF_lastaction{3}\fi% on était dans un cycle
							\CF_in_cyclefalse
							\aftergroup\CF_chemfig_d
							\def\CF_nextaction{\CF_earg\CF_chemfig_c{\CF_molinparen}}%
					\fi
				}%
				{%
					\CF_ifx\CF_molecule\empty
					{%
						\let\CF_nextaction\endgroup
					}
					{% ce qui reste après le noeud courant n'est pas vide, ne commence pas par "-=~", ni par une parenthèse
						\CF_ifx*\CF_toksa% un cycle ?
						{%
							\ifnum\CF_lastaction=2
								\def\CF_currentfromatom{1}% Bugfix 1.6d
								\CF_insertemptygroup\CF_molecule% insère un groupe vide
							\else
								\ifCF_in_cycle
									\def\CF_lastaction{3}%
								\fi% on était dans un cycle
								\ifnum\CF_lastaction=3
									\let\CF_lastcyclenum\CF_cyclenum
								\fi
								\CF_eesecond\CF_iffirsttokis{\expandafter\CF_gobarg\CF_molecule}*%
								{%
									\CF_eesecond{\def\CF_molecule}{\expandafter\CF_gobarg\CF_molecule}% enlève la 1er étoile
									\CF_eesecond\CF_iffirsttokis{\expandafter\CF_gobarg\CF_molecule}[% un crochet ensuite ?
									{%
										\expandafter\CF_parsecyclepreamblewithoptarg\CF_molecule\_nil% \begingroup inclus
									}%
									{%
										\def\CF_cycle_arcstartangle{0}%
										\def\CF_cycle_arcendangle{360}%
										\let\CF_cycle_arcdirecttikz\empty
										\expandafter\CF_parsecyclepreamble\CF_molecule\_nil% \begingroup inclus
									}%
									\CF_cycle_arctrue
								}%
								{%
									\expandafter\CF_parsecyclepreamble\CF_molecule\_nil% \begingroup inclus
									\CF_cycle_arcfalse
								}%
								\CF_cnt_cyclebonds0
								\edef\CF_hookcycle{\CF_bondoutnode}%
								\let\CF_hookatomgroup\CF_previousatomgroup
								\CF_ifnum{\CF_previousatomgroup_nodim=0 }% bugfix v1.71
								{%
									\def\CF_cyclejoinlast{1}% joindre le dernier
								}%
								{%
									\def\CF_cyclejoinlast{0}%
								}%
								\CF_in_cycletrue
								\global\advance\CF_cnt_cycle1
								\CF_ifnum{\CF_lastaction=3 }
								{%
									\pgfmathsetmacro\CF_initcycleangle{360/\CF_lastcyclenum-180}% c'est un cycle dans un cycle
								}
								{%
									\pgfmathsetmacro\CF_initcycleangle{-180/\CF_cyclenum-90+\CF_cycleanglecorrection}%
								}
								\aftergroup\CF_chemfig_d
								\def\CF_nextaction{\CF_earg\CF_chemfig_c{\CF_molinparen}}%
							\fi
						}%
						{%
							\CF_error{something went wrong here: \detokenize\expandafter{\CF_molecule}^^JIf you think it's a bug, please, send a Minimal Example to the author}%
						}%
					}%
				}%
			}%
		}%
			{%
			\expanded{\noexpand\CF_drawatomgroup{\unexpanded\expandafter{\CF_currentangle}}{\unexpanded\expandafter{\CF_currenttoatom}}{\unexpanded\expandafter{\CF_currentatomgroup}}}%
			}%
		}%
	\CF_nextaction
}

\def\CF_parsecyclepreamble*#1#2\_nil{%
	\ifnum#1<3
		\CF_error{a cycle must be at least a triangle.^^JThe number following "*" must be 3 or more}%
	\fi
	\def\CF_molecule{#2}%
	\CF_earg\CF_grabsubmol{\CF_molecule}%
	\begingroup
	\def\CF_cyclenum{#1}%
}

\def\CF_parsecyclepreamblewithoptarg*[#1]#2#3\_nil{%
	\CF_cycleparseoptarg#1,\empty,\empty,\empty\_nil
	\CF_parsecyclepreamble*#2#3\_nil
}

\def\CF_cycleparseoptarg#1,#2,#3\_nil{%
	\CF_ifempty{#1}
	{%
		\def\CF_cycle_arcstartangle{0}%
	}
	{%
		\def\CF_cycle_arcstartangle{#1}%
	}%
	\CF_ifempty{#2}
	{%
		\def\CF_cycle_arcendangle{360}%
	}
	{%
		\def\CF_cycle_arcendangle{#2}%
	}%
	\expandafter\def\expandafter\CF_cycle_arcdirecttikz\expandafter{\CF_sanitizelastitem#3,\empty\_nil}%
}

\def\CF_grabsubmol#1{%
	\begingroup
		\catcode`(1 \catcode`)2
		\expandafter\expandafter\expandafter
	\endgroup
	\expandafter\CF_grabsubmol_a\scantokens{\relax#1}%
}

\def\CF_grabsubmol_a#1\_nil{%
	\expandafter\CF_assigntonil\expandafter\CF_molecule\scantokens\expandafter\expandafter\expandafter{\expandafter\CF_gobarg    \CF_gobarg#1}%
	\expandafter\CF_assigntonil\expandafter\CF_molinparen\scantokens\expandafter\expandafter\expandafter{\expandafter\CF_firsttonil\CF_gobarg#1\_nil}%
}

\def\CF_ifupletter#1{%
	\ifcat\relax\noexpand#1%
		\let\CF_next\CF_execsecond% faux si c'est une sc
	\else
		\ifnum`#1<`A
			\let\CF_next\CF_execsecond
		\else
			\ifnum`#1>`Z
				\let\CF_next\CF_execsecond
			\else
				\let\CF_next\CF_execfirst
			\fi
		\fi
	\fi
	\CF_next
}

% Créé 4 noeuds au dessus et au dessous des noeuds #1 et #2
% à une distance de #3 du noeud #1 et #4 du noeud #2
\def\CF_createnormnodes#1#2#3#4{%
	\CF_doifnotempty{#3}
	{%
		\CF_save
		{%
			\noexpand\node[shape=coordinate,at=(#1),xshift=#3*\CF_normx,yshift=#3*\CF_normy](#11){};^^J%
			\noexpand\node[shape=coordinate,at=(#1),xshift=-#3*\CF_normx,yshift=-#3*\CF_normy](#12){};%
		}%
		\node[shape=coordinate,at=(#1),xshift=#3*\CF_normx,yshift=#3*\CF_normy](#11){};%
		\node[shape=coordinate,at=(#1),xshift=-#3*\CF_normx,yshift=-#3*\CF_normy](#12){};%
	}%
	\CF_doifnotempty{#4}
	{%
		\CF_save
		{%
			\noexpand\node[shape=coordinate,at=(#2),xshift=#4*\CF_normx,yshift=#4*\CF_normy](#21){};^^J%
			\noexpand\node[shape=coordinate,at=(#2),xshift=-#4*\CF_normx,yshift=-#4*\CF_normy](#22){};%
		}%
		\node[shape=coordinate,at=(#2),xshift=#4*\CF_normx,yshift=#4*\CF_normy](#21){};%
		\node[shape=coordinate,at=(#2),xshift=-#4*\CF_normx,yshift=-#4*\CF_normy](#22){};%
	}%
}

\def\CF_extractxy#1#2#3#4{% #1=noeud  #2= ancre  #3=macro pour x  #4=macro pour y
	\pgfextractx\CF_dim_tmp{\pgfpointanchor{#1}{#2}}\edef#3{\the\CF_dim_tmp}%
	\pgfextracty\CF_dim_tmp{\pgfpointanchor{#1}{#2}}\edef#4{\the\CF_dim_tmp}%
}

\def\CF_extractx#1#2#3{%
	\pgfextractx\CF_dim_tmp{\pgfpointanchor{#1}{#2}}\edef#3{\the\CF_dim_tmp}%
}

\def\CF_extracty#1#2#3{%
	\pgfextracty\CF_dim_tmp{\pgfpointanchor{#1}{#2}}\edef#3{\the\CF_dim_tmp}%
}

\def\CF_distancebetweenpoints#1#2#3#4#5{%
	\CF_extractxy{#1}{#2}\CF_dim_tmpax\CF_dim_tmpay
	\CF_extractxy{#3}{#4}\CF_dim_tmpbx\CF_dim_tmpby
	\pgfmathsetmacro#5{veclen(\CF_dim_tmpbx-\CF_dim_tmpax,\CF_dim_tmpby-\CF_dim_tmpay)}%
}

\def\CF_computenodevect#1#2{%
	\CF_distancebetweenpoints{#1}{center}{#2}{center}\CF_vectorlen
	\pgfmathsetmacro\CF_normx{(\CF_dim_tmpay-\CF_dim_tmpby)/\CF_vectorlen}%
	\pgfmathsetmacro\CF_normy{(\CF_dim_tmpbx-\CF_dim_tmpax)/\CF_vectorlen}%
}

\def\CF_ifzerodistance#1#2#3#4{% #1=nom du noeud 1  #2=ancre du noeud 1  #3=nom noeud 2  #4=ancre noeud 2  T-> si distance nulle  F->sinon
	\CF_extractxy{#3}{#4}\CF_dim_tmpax\CF_dim_tmpay
	\CF_extractxy{#1}{#2}\CF_dim_tmpbx\CF_dim_tmpby
	\CF_ifnum{0\ifdim\CF_dim_tmpax=\CF_dim_tmpbx\space1\fi\ifdim\CF_dim_tmpay=\CF_dim_tmpby\space1\fi=11 }
}

\def\CF_setoffset#1#2{% #1=macro recevant le résultat  #2= nom du noeud   bugfix 1.71
	\CF_doifempty#1%
	{%
		\CF_ifzerodistance{#2}{south west}{#2}{north east}% si noeud de dimension nulle
		{%
			\def#1{0pt}%
		}
		{%
			\edef#1{\CF_bondoffset}%
		}%
	}%
}

\def\CF_drawbond#1#2#3#4#5{% #1=type de liaison #2 et #3:nom de noeuds de début et fin #4 et #5: sc contenant les atomes de débutet fin
	\CF_setoffset\CF_startoffset{#2}%
	\CF_setoffset\CF_endoffset{#3}%
	\let\CF_currentbondstyle\CF_bondstyle
	\CF_doifnotempty\CF_currenttikz
	{%
		\CF_eaddtomacro\CF_currentbondstyle{\expandafter,\CF_currenttikz}%
	}%
	\CF_save{\noexpand\path(#2)--(#3)coordinate[pos=0](#2@)coordinate[pos=1](#3@);}%
	\path(#2)--(#3)coordinate[pos=0](#2@)coordinate[pos=1](#3@);%
	\CF_computenodevect{#2@}{#3@}%
	\pgfmathsetmacro\CF_startcoeff{\CF_startoffset/\CF_vectorlen}%
	\pgfmathsetmacro\CF_endcoeff{1-\CF_endoffset/\CF_vectorlen}%
	\CF_save{\noexpand\path(#2@)--(#3@)coordinate[pos=\CF_startcoeff](#2@@)coordinate[pos=\CF_endcoeff](#3@@);}%
	\path(#2@)--(#3@)coordinate[pos=\CF_startcoeff](#2@@)coordinate[pos=\CF_endcoeff](#3@@);%
	\CF_doifnotempty\CF_movebondname% on doit poser un noeud sur la liaison
	{%
		\CF_save{\noexpand\path(#2@@)--(#3@@)coordinate[overlay,pos=\CF_movebondcoeff](\CF_movebondname);}%
		\path(#2@@)--(#3@@)coordinate[overlay,pos=\CF_movebondcoeff](\CF_movebondname);%
		\let\CF_movebondname\empty
	}%
	\ifcase#1\relax
		\CF_error{unknown bond type, this error should not occur^^JIf you think it's a bug, send a Minimal Example to the author}%
	\or% 1 = liaison simple
		\ifboolKV[chemfig]{bond join}% nouveau v1.66
		{%
			\scope
				\expandafter\tikzset\expandafter{\CF_currentbondstyle}%
				\ifnum\CF_previousbondtype=5
					\ifnum0\ifdim\CF_startoffset=0pt 1\fi\ifdim\CF_previous_offset=0pt 1\fi=11
						\CF_ifnum{\CF_previousatomgroup_nodim=1 }% bugfix 1.71
						{%
							\CF_dim_turnangle=\CF_compute_turnangle{\CF_previousbondangle}{\CF_currentangle}%
							\pgfmathsetmacro\CF_cram_overlap{\CF_crambasewidth*cos(\CF_dim_turnangle)-\pgflinewidth}%
							\ifdim\CF_cram_overlap pt>0pt
								\CF_warning{Cram base too wide to join}%
							\else
								\CF_save{%
									\ifdim\CF_dim_turnangle>0pt % si on tourne à gauche
										\noexpand\fill(#2@@)--(\CF_previousnode@@2)--([shift=(\CF_currentangle-90:\pgflinewidth/2)]#2@@)--cycle;
									\else
										\noexpand\fill(#2@@)--(\CF_previousnode@@1)--([shift=(\CF_currentangle+90:\pgflinewidth/2)]#2@@)--cycle;
									\fi
								}%
								\ifdim\CF_dim_turnangle>0pt % si on tourne à gauche
									\fill(#2@@)--(\CF_previousnode@@2)--([shift=(\CF_currentangle-90:\pgflinewidth/2)]#2@@)--cycle;
								\else
									\fill(#2@@)--(\CF_previousnode@@1)--([shift=(\CF_currentangle+90:\pgflinewidth/2)]#2@@)--cycle;
								\fi
							\fi
						}
						{}%
					\fi
				\fi
				\expandafter
			\endscope
			\expanded{\def\noexpand\CF_simplebondwidth{\the\dimexpr\pgflinewidth}}% récupère l'épaisseur de ligne si une liaison ">" est à suivre
		}
		{}%
		\CF_drawaxisbond{#2}{#3}% trace la liaison simple dans l'axe
	\or% 2 = liaison double
		\ifCF_in_cycle
			\ifnum\CF_doublebondtype=0
				\def\CF_doublebondtype{1}%
			\fi
			\ifnum\CF_flipstate>0
				\def\CF_doublebondtype{2}%
			\fi
			\pgfmathsetmacro\CF_doublebondlengthcorrection{\CF_doublesep*tan(180/\CF_cyclenum)}%
		\fi
		\ifcase\CF_doublebondtype
			\CF_createnormnodes{#2@@}{#3@@}{\CF_doublesep/2}{\CF_doublesep/2}%
			\CF_drawbond_a(#2@@1)--(#3@@1);
			\CF_drawbond_a(#2@@2)--(#3@@2);
			\let\CF_joinbond\CF_zero
		\or
			\CF_createnormnodes{#2@@}{#3@@}\CF_doublesep\CF_doublesep
			\CF_drawaxisbond{#2}{#3}% trace la liaison simple dans l'axe\CF_drawbond_a(#2@@)--(#3@@);
			\begingroup% ajuste éventuellement les longueurs des liaisons doubles
				\ifCF_in_cycle
					\ifdim\CF_startoffset=0pt
						\CF_xaddtomacro\CF_currentbondstyle{,shorten <=\CF_doublebondlengthcorrection pt}%
					\fi
					\ifdim\CF_endoffset=0pt
						\CF_xaddtomacro\CF_currentbondstyle{,shorten >=\CF_doublebondlengthcorrection pt}%
					\fi
				\fi
				\CF_drawbond_a(#2@@1)--(#3@@1);
			\endgroup
		\or
			\CF_createnormnodes{#2@@}{#3@@}\CF_doublesep\CF_doublesep
			\CF_drawaxisbond{#2}{#3}% trace la liaison simple dans l'axe\CF_drawbond_a(#2@@)--(#3@@);
			\begingroup% ajuste éventuellement les longueurs des liaisons doubles
				\ifCF_in_cycle
					\ifdim\CF_startoffset=0pt
						\CF_xaddtomacro\CF_currentbondstyle{,shorten \ifnum\CF_flipstate=0 <=-\else>=\fi\CF_doublebondlengthcorrection pt}%
					\fi
					\ifdim\CF_endoffset=0pt
						\CF_xaddtomacro\CF_currentbondstyle{,shorten \ifnum\CF_flipstate=0 >=-\else<=\fi\CF_doublebondlengthcorrection pt}%
					\fi
				\fi
				\CF_drawbond_a(#2@@2)--(#3@@2);
			\endgroup
		\fi
	\or% 3 = liaison triple
		\CF_createnormnodes{#2@@}{#3@@}\CF_doublesep\CF_doublesep
		\CF_drawaxisbond{#2}{#3}% trace la liaison simple dans l'axe\CF_drawbond_a(#2@@)--(#3@@);
		\CF_drawbond_a(#2@@1)--(#3@@1);
		\CF_drawbond_a(#2@@2)--(#3@@2);
	\or% 4 = liaison Cram pleine de #2 vers #3 : >
		\CF_createnormnodes{#2@@}{#3@@}{\CF_crambasewidth/2}{\pgflinewidth/2}%
		\CF_ifnum{\CF_previousatomgroup_nodim=1 }% bugfix 1.71
		{%
			\ifnum\ifboolKV[chemfig]{bond join}10\ifdim\CF_startoffset=0pt 1\fi\ifdim\CF_previous_offset=0pt 1\fi=111
				\ifnum\CF_previousbondtype=1
					\CF_dim_turnangle=\CF_compute_turnangle{\CF_previousbondangle}{\CF_currentangle}%
					\pgfmathsetmacro\CF_cram_overlap{\CF_crambasewidth*cos(\CF_dim_turnangle)-\CF_simplebondwidth}%
					\ifdim\CF_cram_overlap pt>0pt
						\CF_warning{Cram base too wide to join}%
					\else
						\CF_save{\CF_eafter{\noexpand\fill[}\CF_currentbondstyle](#2@@1)--([shift=(\CF_previousbondangle\ifdim\CF_dim_turnangle>0pt -\else+\fi90:\CF_simplebondwidth/2)]\CF_previousnode@@)--(#2@@2)--cycle;}%
						\CF_eafter{\fill[}\CF_currentbondstyle](#2@@1)--([shift=(\CF_previousbondangle\ifdim\CF_dim_turnangle>0pt -\else+\fi90:\CF_simplebondwidth/2)]\CF_previousnode@@)--(#2@@2)--cycle;%
					\fi
				\else
					\ifnum\CF_previousbondtype=5
						\CF_dim_turnangle=\CF_compute_turnangle{\CF_previousbondangle}{\CF_currentangle}%
						\CF_save{\CF_eafter{\noexpand\fill[}\CF_currentbondstyle](#2@@1)--([shift=(\CF_previousbondangle\ifdim\CF_dim_turnangle>0pt -\else+\fi90:\CF_crambasewidth/2)]\CF_previousnode@@)--(#2@@2)--cycle;}%
						\CF_eafter{\fill[}\CF_currentbondstyle](#2@@1)--([shift=(\CF_previousbondangle\ifdim\CF_dim_turnangle>0pt -\else+\fi90:\CF_crambasewidth/2)]\CF_previousnode@@)--(#2@@2)--cycle;%
					\fi
				\fi
			\fi
		}
		{}%
		\ifboolKV[chemfig]{bond join}
			{\CF_eafter{\fill[}}% pas de contour si liaisons jointes
			{\CF_eafter{\filldraw[}}
			\CF_currentbondstyle](#2@@1)--(#2@@2)--(#3@@2)--(#3@@1)--cycle;%
		\CF_save{%
			\ifboolKV[chemfig]{bond join}
				{\CF_eafter{\noexpand\fill[}}% pas de contour si liaisons jointes
				{\CF_eafter{\noexpand\filldraw[}}
				\CF_currentbondstyle](#2@@1)--(#2@@2)--(#3@@2)--(#3@@1)--cycle;%
		}%
		\let\CF_joinbond\CF_zero
	\or% 5 = liaison Cram pleine de #2 vers #3 : <
		\CF_createnormnodes{#3@@}{#2@@}{\CF_crambasewidth/2}{\pgflinewidth/2}%
		\CF_save{%
			\ifboolKV[chemfig]{bond join}
				{\CF_eafter{\noexpand\fill[}}% pas de contour si liaisons jointes
				{\CF_eafter{\noexpand\filldraw[}}
				\CF_currentbondstyle](#2@@2)--(#2@@1)--(#3@@1)--(#3@@2)--cycle;
		}%
		\ifboolKV[chemfig]{bond join}
			{\CF_eafter{\fill[}}% pas de contour si liaisons jointes
			{\CF_eafter{\filldraw[}}
			\CF_currentbondstyle](#2@@2)--(#2@@1)--(#3@@1)--(#3@@2)--cycle;
		\let\CF_joinbond\CF_zero
	\or% 6 = liaison Cram pointillée de #2 vers #3
		\scope
			\CF_createnormnodes{#2@@}{#3@@}{\CF_crambasewidth/2}{}%
			\CF_clipcramornot(#2@@1)--(#2@@2)--(#3@@)--(#2@@1);
			\CF_save{\CF_eafter{\noexpand\draw[}\CF_currentbondstyle,dash pattern=on \CF_cramdashlength off \CF_cramdashsep,line width=\CF_crambasewidth](#2@@)--(#3@@);}%
			\CF_eafter{\draw[}\CF_currentbondstyle,dash pattern=on \CF_cramdashlength off \CF_cramdashsep,line width=\CF_crambasewidth](#2@@)--(#3@@);
		\endscope
		\let\CF_joinbond\CF_zero
	\or% 7 = liaison Cram pointillée de #3 vers #2
		\scope
			\CF_createnormnodes{#3@@}{#2@@}{\CF_crambasewidth/2}{}%
			\CF_clipcramornot(#3@@1)--(#3@@2)--(#2@@)--(#3@@1);
			\CF_save{\CF_eafter{\noexpand\draw[}\CF_currentbondstyle,dash pattern=on \CF_cramdashlength off \CF_cramdashsep,line width=\CF_crambasewidth](#3@@)--(#2@@);}%
			\CF_eafter{\draw[}\CF_currentbondstyle,dash pattern=on \CF_cramdashlength off \CF_cramdashsep,line width=\CF_crambasewidth](#3@@)--(#2@@);
		\endscope
		\let\CF_joinbond\CF_zero
	\or% 8 = liaison cram rectangle évidé de #2 vers #3
		\CF_createnormnodes{#2@@}{#3@@}{\CF_crambasewidth/2}{}%
		\CF_save{\CF_eafter{\noexpand\draw[}\CF_currentbondstyle,line join=bevel](#2@@1)--(#2@@2)--(#3@@)--cycle;}%
		\CF_eafter{\draw[}\CF_currentbondstyle,line join=bevel](#2@@1)--(#2@@2)--(#3@@)--cycle;
		\let\CF_joinbond\CF_zero
	\or% 9 = liaison cram rectangle évidé de #3 vers #1
		\CF_createnormnodes{#3@@}{#2@@}{\CF_crambasewidth/2}{}%
		\CF_save{\CF_eafter{\noexpand\draw[}\CF_currentbondstyle,line join=bevel](#3@@1)--(#3@@2)--(#2@@)--cycle;}%
		\CF_eafter{\draw[}\CF_currentbondstyle,line join=bevel](#3@@1)--(#3@@2)--(#2@@)--cycle;
		\let\CF_joinbond\CF_zero
	\else
		\CF_error{unknown bond type, this error should not occur^^JIf you think it's a bug, send a Minimal Example to the author}%
	\fi
	\let\CF_startoffset\empty
	\let\CF_previous_offset\CF_endoffset
	\let\CF_endoffset\empty
	\let\CF_previoustikz\CF_currenttikz
	\let\CF_previousbondangle\CF_previousangle
	\edef\CF_previousnode{#3}%
}

\def\CF_compute_turnangle#1#2{%
	\expandafter\CF_compute_turnangle_a\the\dimexpr #2pt-#1pt\relax\_nil
}
\def\CF_compute_turnangle_a#1\_nil{%
	\dimexpr#1
		\ifdim#1<-180pt +360pt\else
		\ifdim#1>180pt  -360pt\fi\fi\relax
}

\def\CF_drawaxisbondnojoin#1#2{\CF_drawbond_a(#1@@)--(#2@@);}

\def\CF_drawaxisbondjoin#1#2{% dessine une liaison simple dans l'axe avec raccord rétrograde
	\ifCF_in_cycle\ifnum\CF_cnt_cyclebonds=\CF_cyclenum\relax
		\let\CF_nexttikz\CF_cyclefirsttikz
	\fi\fi
	\ifnum\CF_joinbond=0
		\ifCF_in_cycle
			\ifnum\CF_cnt_cyclebonds=\CF_cyclenum\relax
				\ifnum\CF_cyclejoinlast=1
					\CF_drawbond_a(#1@@)--(#2@@)--%
						([shift=(\CF_previousbondangle+2*\CF_cycleincrementangle:.5\pgflinewidth)]#2@@);
				\else
					\CF_drawbond_a(#1@@)--(#2@@);
				\fi
			\else
				\CF_drawbond_a(#1@@)--(#2@@);
			\fi
		\else
			\CF_drawbond_a(#1@@)--(#2@@);
		\fi
		\def\CF_joinbond{1}%
	\else
		\CF_ifx\CF_previoustikz\CF_currenttikz
		{%
			\def\CF_joinbond{1}%
			\CF_ifnum{\CF_previousatomgroup_nodim=1 }% bugfix 1.71
			{%
				\CF_ifx\CF_previousbondangle\empty% si début molécule
				{%
					\CF_drawbond_a(#1@@)--(#2@@);
				}
				{%
					\ifdim\CF_startoffset=0pt
						\ifCF_in_cycle
							\ifnum\CF_cnt_cyclebonds=\CF_cyclenum\relax
								\ifnum\CF_cyclejoinlast=1
									\CF_ifx\CF_cyclefirsttikz\CF_currenttikz
									{%
										\CF_drawbond_a([shift=(\CF_previousbondangle:-.5\pgflinewidth)]#1@@)--(#1@@)--(#2@@)--([shift=(\CF_previousbondangle+2*\CF_cycleincrementangle:.5\pgflinewidth)]#2@@);
									}
									{%
										\CF_drawbond_a([shift=(\CF_previousbondangle:-.5\pgflinewidth)]#1@@)--(#1@@)--(#2@@);
									}%
								\else
									\CF_drawbond_a([shift=(\CF_previousbondangle:-.5\pgflinewidth)]#1@@)--(#1@@)--(#2@@);
								\fi
							\else
								\CF_drawbond_a([shift=(\CF_previousbondangle:-.5\pgflinewidth)]#1@@)--(#1@@)--(#2@@);
							\fi
						\else
							\CF_drawbond_a([shift=(\CF_previousbondangle:-.5\pgflinewidth)]#1@@)--(#1@@)--(#2@@);
						\fi
					\else
						\CF_drawbond_a(#1@@)--(#2@@);
					\fi
				}%
			}
			{%
				\CF_drawbond_a(#1@@)--(#2@@);%
			}%
		}
		{%
			\ifCF_in_cycle
				\ifnum\CF_cnt_cyclebonds=\CF_cyclenum\relax
					\ifnum\CF_cyclejoinlast=1
						\CF_ifx\CF_nexttikz\CF_currenttikz
						{%
							\CF_drawbond_a(#1@@)--(#2@@)--([shift=(\CF_previousbondangle+2*\CF_cycleincrementangle:.5\pgflinewidth)]#2@@);%
						}%
						{%
							\CF_drawbond_a(#1@@)--(#2@@);%
						}%
					\else
						\CF_drawbond_a(#1@@)--(#2@@);
					\fi
				\else
					\CF_drawbond_a(#1@@)--(#2@@);
				\fi
			\else
				\CF_drawbond_a(#1@@)--(#2@@);%
			\fi
		}%
	\fi
	\ifdim\CF_endoffset=0pt \else
		\let\CF_joinbond\CF_zero
	\fi
}

\def\CF_drawbond_a#1;{%
	\CF_save{\CF_eafter{\noexpand\draw[}\CF_currentbondstyle]#1;}
	\CF_eafter{\draw[}\CF_currentbondstyle]#1;%
}

\def\CF_drawallhooks{% dessine tous les crochets contenus dans la sc \CF_hookdraw
	\CF_doifnotempty\CF_hookdrawlist
	{%
		\expandafter\CF_drawfirsthook\CF_hookdrawlist\_nil% trace un lien de crochet à crochet
		\CF_drawallhooks
	}%
}

\def\CF_drawfirsthook[#1,#2,#3]#4#5#6#7#8\_nil{%
	\def\CF_hookdrawlist{#8}%
	\begingroup
		\let\CF_joinbond\CF_zero
		\def\CF_currenttikz{#3}%
		\def\CF_hookstartcontent{#6}%
		\def\CF_hookendcontent{#7}%
		\CF_ifinteger{#2}%
		{%
			\CF_drawbond{#2}{#4}{#5}\CF_hookstartcontent\CF_hookendcontent
		}%
		{%
			\CF_assignbondcode{#2}\CF_bondcurrentnum
			\CF_drawbond\CF_bondcurrentnum{#4}{#5}\CF_hookstartcontent\CF_hookendcontent
		}%
	\endgroup
}

\def\CF_extractatom#1-#2\_nil{#2}% transforme le bound@outnode en n° de l'atome

\def\CF_gobblemovearg @#1#2\_nil#3{%
	\expandafter\def\csname atom_\number\CF_cnt_atom\endcsname{#2}%
	\CF_ifinstr{#1},%
	{%
		\CF_analysemovearg#1\_nil#3\let\CF_movebondcoeff\empty
	}%
	{%
		\def#3{#1}%
	}%
	\CF_doifempty{#2}
	{%
		\let\CF_nodestrut\empty
	}%
}%

\def\hflipnext{\def\CF_flipstate{1}}

\def\vflipnext{\def\CF_flipstate{2}}
\let\CF_flipstate\CF_zero

\def\CF_drawatomgroup#1#2#3{% #1=angle d'arrivée de la liaison #2=numero atome sur lequel arrive la liaison #3=groupe d'atomes
	\expandafter\let\expandafter\CF_bondoutcontent% assigne le contenu de l'atome d'où part la liaison
		\csname
			\ifdefined\CF_bondoutnode
				atom_\expandafter\CF_extractatom\CF_bondoutnode\_nil
			\else
				empty%
			\fi
		\endcsname
	\global\advance\CF_cnt_group1
	\let\CF_currentatom\empty
	\global\let\CF_hookdrawlist\empty
	\CF_cnt_atomgroup0 % est le nombre d'atome dans le groupe que va calculer \CF_drawatomgroup_a
	\CF_iffirsttokis{#3}?%
	{%
		\CF_drawatomgroup_a{{}#3}%
	}
	{%
		\CF_drawatomgroup_a{#3}%
	}%
	\def\CF_currentatomgroup{#3}%
	\CF_removemovearg\CF_currentatomgroup% enlève les "@{<nom>}"
	\CF_ifinstr{#3}?%
	{%
		\CF_removehook\CF_currentatomgroup
		\CF_earg\CF_stripsp\CF_currentatomgroup\CF_ifempty
		{%
			\let\CF_currentatomgroup\empty% rollback à 1.66
		}
		{}%
	}%
	{}%
	\CF_doifnotempty{#2}
		{%
		\ifnum#2<1
			\CF_warning{no atom found at position #2, pershaps you mispelled the optional argument of the bond.}%
		\else
			\ifnum#2>\CF_cnt_atomgroup
				\CF_error{no atom found at position #2, pershaps you mispelled the optional argument of the bond.}%
			\fi
		\fi
		}%
	\edef\CF_hookatomnumber{%
		\CF_ifempty{#2}
		{%
			\ifdim#1pt>90pt
				\ifdim#1pt<270pt
					\number\CF_cnt_atomgroup
				\else
					1%
				\fi
			\else
				1%
			\fi
		}
		{%
			#2%
		}%
	}%
	\CF_cnt_atom\CF_hookatomnumber
	\ifboolKV[chemfig]{use atom strut}% bugfix 1.71
	{% branche inutile si use atom strut=false
		\CF_ifzerodim\CF_currentatomgroup
		{%
			\let\CF_nodestrut\empty
		}
		{%
			\CF_ifx\empty\CF_bondoutcontentsaved
			{%
				\def\CF_nodestrut{\vphantom\CF_bondoutcontent}%
			}%
			{%
				\def\CF_nodestrut{\vphantom\CF_bondoutcontentsaved}%
			}%
		}%
	}
	{%
		\let\CF_nodestrut\empty
	}%
	\edef\CF_optstring{anchor=\ifnum\CF_lastaction=0 base\else\ifCF_in_cycle center\else\ifCF_macro_fixed_bond_length 180+#1\else center\fi\fi\fi,at=(CF_node),\CF_nodestyle}% premier atome de la molécule affiché
	\loop
		\unless\ifnum\CF_cnt_atom>\CF_cnt_atomgroup
		\CF_eesecond\CF_iffirsttokis{\csname atom_\number\CF_cnt_atom\endcsname}@% l'atome courant commence par un "@"
		{%
			\expandafter\expandafter\expandafter\CF_gobblemovearg\csname atom_\number\CF_cnt_atom\endcsname\_nil\CF_moveatomname
			\edef\CF_nodecontent{\noexpand\printatom{\unexpanded\expandafter\expandafter\expandafter{\csname atom_\number\CF_cnt_atom\endcsname}\unexpanded\expandafter{\CF_nodestrut}}}%
			\CF_eafter{\node[}\CF_optstring,overlay](\CF_moveatomname){\phantom{\CF_nodecontent}};%
			\CF_save{\CF_eafter{\noexpand\node[}\CF_optstring,overlay](\CF_moveatomname){\noexpand\phantom{\unexpanded\expandafter{\CF_nodecontent}}};}%
			\let\CF_moveatomname\empty
		}
		{%
			\edef\CF_nodecontent{\noexpand\printatom{\unexpanded\expandafter\expandafter\expandafter{\csname atom_\number\CF_cnt_atom\endcsname}\unexpanded\expandafter{\CF_nodestrut}}}%
		}%
		\ifboolKV[chemfig]{debug}
		{%
			\CF_eafter{\node[}\CF_optstring,draw=gray](n\number\CF_cnt_group-\number\CF_cnt_atom){\CF_nodecontent};%
			\CF_save{\CF_eafter{\noexpand\node[}\CF_optstring,draw=gray](n\number\CF_cnt_group-\number\CF_cnt_atom){\unexpanded\expandafter{\CF_nodecontent}};}%
			\CF_show_debug_atom
		}
		{%
			\CF_eafter{\node[}\CF_optstring](n\number\CF_cnt_group-\number\CF_cnt_atom){\CF_nodecontent};%
			\CF_save{\CF_eafter{\noexpand\node[}\CF_optstring](n\number\CF_cnt_group-\number\CF_cnt_atom){\unexpanded\expandafter{\CF_nodecontent}};}%
		}%
		\let\CF_nodestrut\empty
		\advance\CF_cnt_atom1
		\edef\CF_optstring{anchor=base \ifnum\CF_flipstate=1 east\else west\fi,at=(n\number\CF_cnt_group-\number\numexpr\CF_cnt_atom-1.base \ifnum\CF_flipstate=1 west\else east\fi),\CF_nodestyle}%
	\repeat
	\CF_cnt_atom\CF_hookatomnumber
	\ifnum\CF_lastaction=2 % s'il faut tracer une liaison
		\gdef\CF_cycleanglecorrection{0}% alors c'est qu'un cycle ne peut pas commencer la molécule : annulation de la correction d'angle
		\CF_drawbond\CF_bondtype{\CF_bondoutnode}{n\number\CF_cnt_group-\number\CF_cnt_atom}\CF_previousatomgroup\CF_currentatomgroup
		\let\CF_previousbondtype\CF_bondtype
	\fi
	\def\CF_lastaction{1}% met la dernière action à 1 : affichage d'un noeud
	\loop
		\ifnum\CF_cnt_atom>1
		\advance\CF_cnt_atom-1
		\edef\CF_optstring{anchor=base \ifnum\CF_flipstate=1 west\else east\fi,at=(n\number\CF_cnt_group-\number\numexpr\CF_cnt_atom+1.base \ifnum\CF_flipstate=1 east\else west\fi),\CF_nodestyle}%
		\CF_eesecond\CF_iffirsttokis{\csname atom_\number\CF_cnt_atom\endcsname}@% l'atome courant commence par un "@"
		{%
			\expandafter\expandafter\expandafter\CF_gobblemovearg\csname atom_\number\CF_cnt_atom\endcsname\_nil\CF_moveatomname
			\edef\CF_nodecontent{\noexpand\printatom{\unexpanded\expandafter\expandafter\expandafter{\csname atom_\number\CF_cnt_atom\endcsname}\unexpanded\expandafter{\CF_nodestrut}}}
			\CF_eafter{\node[}\CF_optstring,overlay](\CF_moveatomname){\phantom{\CF_nodecontent}};%
			\CF_save{\CF_eafter{\noexpand\node[}\CF_optstring,overlay](\CF_moveatomname){\noexpand\phantom{\unexpanded\expandafter{\CF_nodecontent}}};}%
			\let\CF_moveatomname\empty
		}
		{%
			\edef\CF_nodecontent{\noexpand\printatom{\unexpanded\expandafter\expandafter\expandafter{\csname atom_\number\CF_cnt_atom\endcsname}\unexpanded\expandafter{\CF_nodestrut}}}
		}%
		\ifboolKV[chemfig]{debug}
		{%
			\CF_eafter{\node[}\CF_optstring,draw=gray](n\number\CF_cnt_group-\number\CF_cnt_atom){\CF_nodecontent};%
			\CF_save{\CF_eafter{\noexpand\node[}\CF_optstring,draw=gray](n\number\CF_cnt_group-\number\CF_cnt_atom){\unexpanded\expandafter{\CF_nodecontent}};}%
			\CF_show_debug_atom
		}
		{%
			\CF_eafter{\node[}\CF_optstring](n\number\CF_cnt_group-\number\CF_cnt_atom){\CF_nodecontent};%
			\CF_save{\CF_eafter{\noexpand\node[}\CF_optstring](n\number\CF_cnt_group-\number\CF_cnt_atom){\unexpanded\expandafter{\CF_nodecontent}};}%
		}%
	\repeat
	\ifboolKV[chemfig]{debug}
		\CF_show_debug_atomgroup
		{}%
	\CF_drawallhooks
	\edef\CF_lastgroupnumber{\number\CF_cnt_group}%
	\let\CF_previousatomgroup\CF_currentatomgroup
	\CF_ifzerodistance{n\number\CF_cnt_group-1}{south west}{n\number\CF_cnt_group-\number\CF_cnt_atomgroup}{north west}% new v1.71
		{\def\CF_previousatomgroup_nodim{1}}
		{\def\CF_previousatomgroup_nodim{0}}%
}

\def\CF_show_debug_atom{%
	\CF_save{\noexpand\node[at=(n\number\CF_cnt_group-\number\CF_cnt_atom.south),anchor=north,outer sep=1pt,overlay]{$\noexpand\scriptscriptstyle\noexpand\color{gray}\number\CF_cnt_atom$};}
	\node[at=(n\number\CF_cnt_group-\number\CF_cnt_atom.south),anchor=north,outer sep=1pt,overlay]{$\scriptscriptstyle\color{gray}\number\CF_cnt_atom$};%
}

\def\CF_show_debug_atomgroup{%
	\CF_save{%
		\noexpand\draw[red,overlay] ([xshift=-.5pt,yshift=.5pt]n\number\CF_cnt_group-1.north west) rectangle ([xshift=.5pt,yshift=-.5pt]n\number\CF_cnt_group-\number\CF_cnt_atomgroup.south east);^^J%
		\noexpand\path (n\number\CF_cnt_group-1.north west) -- (n\number\CF_cnt_group-\number\CF_cnt_atomgroup.north east)%
			node [midway,yshift=1pt,overlay] {$\noexpand\scriptscriptstyle\noexpand\color{red}\number\CF_cnt_group$};%
	}%
	\draw[red,overlay] ([xshift=-.5pt,yshift=.5pt]n\number\CF_cnt_group-1.north west) rectangle ([xshift=.5pt,yshift=-.5pt]n\number\CF_cnt_group-\number\CF_cnt_atomgroup.south east);%
	\path (n\number\CF_cnt_group-1.north west) -- (n\number\CF_cnt_group-\number\CF_cnt_atomgroup.north east)
		node [midway,yshift=1pt,overlay] {$\scriptscriptstyle\color{red}\number\CF_cnt_group$};% TODO : supprimer ce noeud qui ne donne aucune information utile ?
}

\def\CF_savemovearg @#1#2\_nil{\def\CF_currentatom{@{#1}}}

\def\CF_drawatomgroup_a#1{% transforme #1 en un groupe d'atomes
	\CF_ifempty{#1}
	{%
		\expandafter\let\csname atom_\number\CF_cnt_atomgroup\endcsname\CF_currentatom
	}
	{%
		\advance\CF_cnt_atomgroup1
		\CF_iffirsttokis{#1}@%
		{%
			\CF_savemovearg#1\_nil
			\CF_removemovearg_a#1\_nil\CF_aftermovearg
			\CF_earg\CF_drawatomgroup_b{\CF_aftermovearg}%
		}%
		{%
			\let\CF_currentatom\empty
			\CF_drawatomgroup_b{#1}%
		}%
	}%
}

\def\CF_drawatomgroup_b#1{%
	\CF_ifempty{#1}
	{%
		\expandafter\let\csname atom_\number\CF_cnt_atomgroup\endcsname\CF_currentatom
	}
	{%
		\futurelet\CF_toksa\CF_gobtonil#1\_nil
		\CF_ifxcase\CF_toksa
			\bgroup
			{%
				\CF_eaddtomacro\CF_currentatom{\expandafter{\CF_firsttonil#1\_nil}}%
				\CF_earg\CF_drawatomgroupc{\CF_gobarg#1}%
			}%
			\CF_sptoken
			{%
				\CF_addtomacro\CF_currentatom{ }%
				\CF_earg\CF_drawatomgroupc{\CF_afterspace#1\_nil}%
			}%
		\CF_ifxcase_elseif
			\CF_eaddtomacro\CF_currentatom{\CF_firsttonil#1\_nil}%
			\CF_earg\CF_drawatomgroupc{\CF_gobarg#1}%
		\CF_ifxcase_endif
	}%
}

% enlève tous les "@{nom}" de la sc #1
\def\CF_removemovearg#1{%
	\CF_earg\CF_ifinstr{#1}@%
	{%
		\expandafter\CF_removemovearg_a#1\_nil#1%
		\CF_removemovearg#1%
	}%
	{}%
}

% enlève le premier "@{<nom>}" de l'argument et l'assigne à #2
\def\CF_removemovearg_a#1\_nil#2{%
	\def\CF_removemovearg_b##1@{%
		\expandafter\def\expandafter#2\expandafter{\CF_gobarg##1}% mange le \relax
		\CF_removemoveargc\relax
	}%
	\def\CF_removemoveargc##1\_nil{\CF_eaddtomacro#2{\CF_gobtwoargs##1}}% mange le \relax et le <nom>
	\CF_removemovearg_b\relax#1\_nil
}

\def\CF_drawatomgroupc#1{% transforme #1 en un groupe d'atomes
	\CF_ifempty{#1}
	{%
		\expandafter\let\csname atom_\number\CF_cnt_atomgroup\endcsname\CF_currentatom
		\let\CF_currentatom\empty
	}
	{%
		\CF_iffirsttokis{#1}@%
		{%
			\expandafter\let\csname atom_\number\CF_cnt_atomgroup\endcsname\CF_currentatom
			\let\CF_currentatom\empty
			\CF_drawatomgroup_a{#1}%
		}%
		{%
			\CF_ifx|\CF_toksa
			{%
				\expandafter\let\csname atom_\number\CF_cnt_atomgroup\endcsname\CF_currentatom
				\let\CF_currentatom\empty
				\CF_earg\CF_drawatomgroup_a{\CF_gobarg#1}%
			}%
			{%
				\CF_ifx\CF_sptoken\CF_toksa
				{%
					\CF_addtomacro\CF_currentatom{ }%
					\CF_earg\CF_drawatomgroupc{\CF_afterspace#1\_nil}%
				}%
				{%
					\CF_ifx\bgroup\CF_toksa
					{%
						\CF_eaddtomacro\CF_currentatom{\expandafter{\CF_firsttonil#1\_nil}}%
						\CF_earg\CF_drawatomgroupc{\CF_gobarg#1}%
					}%
					{%
						\CF_earg\CF_ifupletter{\CF_firsttonil#1\_nil}%
						{%
							\expandafter\let\csname atom_\number\CF_cnt_atomgroup\endcsname\CF_currentatom
							\let\CF_currentatom\empty
							\CF_drawatomgroup_a{#1}%
						}%
						{%
							\CF_ifx?\CF_toksa
							{%
								\CF_earg\CF_iffirsttokis{\CF_gobarg#1}[% un crochet après le "?"
								{%
									\expandafter\CF_graboptarg\CF_gobarg#1\_nil\CF_afterhook
								}%
								{%
									\CF_eafter{\CF_graboptarg[]}{\CF_gobarg#1}\_nil\CF_afterhook
								}%
								\CF_etwoargs\CF_ifinstr{\CF_hooklist}{\expandafter(\CF_hookcurrentname)}% crochet déjà défini ?
								{%
									\CF_earg\CF_hookparselist{\CF_hookcurrentname}% chercher les caractéristiques du crochet sauvegardé
									\CF_xaddtomacro\CF_hookdrawlist{%
										[\CF_hookcurrentname,\CF_hookcurrentlink,\CF_hookcurrenttikz]{\CF_hooksavedcoord}{n\number\CF_cnt_group-\number\CF_cnt_atomgroup}}%
									\CF_eaddtomacro\CF_hookdrawlist{\expandafter{\CF_hooksavedcontent}}%
									\CF_eaddtomacro\CF_hookdrawlist{\expandafter{\CF_currentatom}}% ajoute les 4 arguments à la liste des crochets à tracer
									\global\let\CF_hookdrawlist\CF_hookdrawlist
								}%
								{%
									\CF_xaddtomacro\CF_hooklist{(\CF_hookcurrentname)|n\number\CF_cnt_group-\number\CF_cnt_atomgroup|}%
									\CF_eaddtomacro\CF_hooklist{\CF_currentatom|}%
									\global\let\CF_hooklist\CF_hooklist
								}%
								\CF_earg\CF_drawatomgroupc{\CF_afterhook}%
							}%
							{%
								\CF_eaddtomacro\CF_currentatom{\CF_firsttonil#1\_nil}%
								\CF_earg\CF_drawatomgroupc{\CF_gobarg#1}%
							}%
						}%
					}%
				}%
			}%
		}%
	}%
}

\def\CF_kookdefaultname{a}
\def\CF_hookdefaultlink{-}
\def\CF_hookdefaulttikz{}

\def\CF_hookparseoptarg#1,#2,#3\_nil{%
	\CF_testemptyandassign\CF_hookcurrentname{#1}\CF_kookdefaultname
	\CF_testemptyandassign\CF_hookcurrentlink{#2}\CF_hookdefaultlink
	\CF_testemptyandassign\CF_hookcurrenttikz{#3}\CF_hookdefaulttikz
}

\def\CF_graboptarg[#1]#2\_nil#3{%
	\CF_hookparseoptarg#1,,\_nil
	\def#3{#2}%
}

\def\CF_hookparselist#1{% #1 est le nom du noeud à retrouver
	\def\CF_hookparselist_a##1(#1)|##2|##3|##4\_nil{%
		\def\CF_hooksavedcoord{##2}%
		\def\CF_hooksavedcontent{##3}%
	}%
	\expandafter\CF_hookparselist_a\CF_hooklist\_nil
}

\def\CF_removehook#1{%
	\CF_earg\CF_ifinstr{#1}?%
	{%
		\CF_eafter{\CF_removehook_a\relax}#1\_nil#1%
		\CF_removehook#1%
	}
	{}%
}

\def\CF_removehook_a#1?#2\_nil#3{%
	\CF_iffirsttokis{#2}[%
	{%
		\CF_removehook_b#1?#2\_nil#3%
	}
	{%
		\expandafter\def\expandafter#3\expandafter{\CF_gobarg#1#2}%
	}%
}

\def\CF_removehook_b#1?[#2]#3\_nil#4{\expandafter\def\expandafter#4\expandafter{\CF_gobarg#1#3}}

\defKV[charge]{%
	.radius      = \CF_defifempty\CF_dotradius {#1}{0.15ex},
	:sep         = \CF_defifempty\CF_dotsep    {#1}{0.3em},
	.style       = \CF_defifempty\CF_dotstyle  {#1}{fill=black},
	"length      = \CF_defifempty\CF_rectlength{#1}{1.5ex},
	"width       = \CF_defifempty\CF_rectwidth {#1}{.3ex}
}
\def\setcharge#{\setKV[charge]}
\def\resetcharge{\restoreKV[charge]}
\setKVdefault[charge]{%
	debug        = false,% trace les contours des noeuds
	macro atom   = \printatom,%macro qui prendra comme argument l'atome recevant la charge
	circle       = false,% false => noeud atome = rectangle
	macro charge = ,% macro attendue (\printatom ou  \ensuremath, par exemple) qui prendra comme argument la charge
	extra sep    = 1.5pt,% séparation additionnelle entre le noeud (cercle ou rectangle) et la position des charges
	overlay      = true,% charges en "surimpression"
	shortcuts    = true,% raccourcis \. \: \| et \" actifs pour Lewis
	lewisautorot = true,% rotation auto charge Lewis
	.radius      = 0.15ex,% rayon du point
	:sep         = 0.3em,% séparation des deux points
	.style       = {fill=black},% style des points
	"length      = 1.5ex,% longueur rectangle
	"width       = .3ex,% largeur rectangle
	"style       = {black,line width=0.4pt},% style rectangle
	|style       = {black,line width=0.4pt},% style ligne
}%
\def\chargedot{\CF_testopt\chargedot_a{}}
\def\chargedot_a[#1]{%
	\begingroup
		\setKV[charge]{#1}%
		\CF_eafter{\tikz\draw[}{\CF_dotstyle}](0,0)circle(\CF_dotradius);%
	\endgroup
}
\def\chargeddot{\CF_testopt\chargeddot_a{}}
\def\chargeddot_a[#1]{%
	\begingroup
		\setKV[charge]{#1}%
		\ifboolKV[charge]{lewisautorot}
		{%
			\pgfmathsetmacro\CF_lewisrot{90+\chargeangle}%
		}
		{%
			\def\CF_lewisrot{0}%
		}%
		\pgfmathsetmacro\CF_halfsep{\CF_dotsep/2}%
		\tikzpicture[anchor=center,rotate=\CF_lewisrot]
			\CF_eafter{\draw[}{\CF_dotstyle}](-\CF_halfsep pt,0)circle(\CF_dotradius)(\CF_halfsep pt,0)circle(\CF_dotradius);%
		\endtikzpicture
	\endgroup
}
\def\chargerect{\CF_testopt\chargerect_a{}}
\def\chargerect_a[#1]{%
	\begingroup
		\setKV[charge]{#1}%
		\ifboolKV[charge]{lewisautorot}
		{%
			\pgfmathsetmacro\CF_lewisrot{90+\chargeangle}%
		}
		{%
			\def\CF_lewisrot{0}%
		}%
		\pgfmathsetmacro\CF_halfwidth{\CF_rectwidth/2}%
		\pgfmathsetmacro\CF_halflength{\CF_rectlength/2}%
		\tikzpicture[anchor=center,rotate=\CF_lewisrot]%
			\CF_eeafter{\draw[}{\useKV[charge]{"style}}](-\CF_halflength pt,-\CF_halfwidth pt)rectangle(\CF_halflength pt,\CF_halfwidth pt);% bugfix 1.51
		\endtikzpicture
	\endgroup
}
\def\chargeline{\CF_testopt\chargeline_a{}}
\def\chargeline_a[#1]{%
	\begingroup
		\setKV[charge]{#1}%
		\ifboolKV[charge]{lewisautorot}
		{%
			\pgfmathsetmacro\CF_lewisrot{90+\chargeangle}%
		}
		{%
			\def\CF_lewisrot{0}%
		}%
		\pgfmathsetmacro\CF_halflength{\CF_rectlength/2}%
		\tikzpicture[anchor=center,rotate=\CF_lewisrot]%
			\CF_eeafter{\draw[}{\useKV[charge]{|style}}](-\CF_halflength pt,0)--(\CF_halflength pt,0);% bugfix 1.51
		\endtikzpicture
	\endgroup
}
\def\CF_enableshortcuts{%
	\let\CF_saveddot \.\let\.\chargedot
	\let\CF_savedddot\:\let\:\chargeddot
	\let\CF_savedrect\"\let\"\chargerect
	\let\CF_savedline\|\let\|\chargeline
	\let\enableshortcuts\relax
	\let\disableshortcuts\CF_disableshortcuts
}
\def\CF_disableshortcuts{%
	\let\.\CF_saveddot
	\let\:\CF_savedddot
	\let\"\CF_savedrect
	\let\|\CF_savedline
	\let\enableshortcuts\CF_enableshortcuts
	\let\disableshortcuts\relax
}
\def\charge{%
	\begingroup
		\catcode`\: 12
		\charge_a{true}%
}
\def\Charge{%
	\begingroup
		\catcode`\: 12
		\charge_a{false}%
}
\def\charge_a#1#2{% #1=TF #2=liste emplacements
		\CF_testopt{\charge_b{#1}}{}#2\_nil
}
\def\charge_b#1[#2]#3\_nil{%
		\charge_c{#1}[#2]{#3}%
}
\def\charge_c#1[#2]#3#4{% #1=TF pour overlay, #2= réglages, #3=liste d'emplacements, #4=atome
		\setcharge{overlay=#1,#2}%
		\setbox\CF_box_charge\hbox{\useKV[charge]{macro atom}{#4}}%
		\CF_ifinsidetikz
		{%
			\pgfinterruptpicture
			\def\CF_atendofcharge{\endpgfinterruptpicture}%
		}
		{%
			\let\CF_atendofcharge\empty
		}%
		\expanded{\noexpand
		\tikzpicture[every node/.style={%
					\ifboolKV[charge]{debug}{draw=red,}{}%
					anchor=base,%
					inner sep=0pt,%
					outer sep=0pt,%
					minimum size=0pt},%
				baseline]}%
			\expanded{\noexpand
			\node[%
				\ifboolKV[charge]{circle}{circle,}{}%
				\ifboolKV[charge]{debug}{draw=green,}{}%
				anchor=base%
				]}%
				(atombox)at(0,0)%
				{\copy\CF_box_charge};% noeud contenant l'atome
			\expanded{\noexpand
			\node[%
				\ifboolKV[charge]{circle}{circle,}{}%
				\ifboolKV[charge]{debug}{draw=blue,}{}%
				anchor=base,%
				inner sep=\useKV[charge]{extra sep},%
				overlay%
				]}%
				(atom)at(0,0){%
						\vrule width0pt height\ht\CF_box_charge  depth\dp\CF_box_charge
						\vrule width\wd\CF_box_charge height\CF_zero depth\CF_zero};% noeud pour placer les charges
			\let\enableshortcuts\relax
			\let\disableshortcuts\relax
			\ifboolKV[charge]{shortcuts}
				\CF_enableshortcuts
				{}% l'atome n'est _PAS_ concerné par les racourcis
			\charge_d#3,\CF_quark=%
		\endtikzpicture
		\CF_atendofcharge
	\endgroup
}
\def\charge_d#1={%
	\CF_ifx\CF_quark{#1}%
	{}
	{%
		\CF_striplastsp{#1}\charge_e=% bugfix 1.54
	}%
}
\def\charge_e#1={%
	\CF_ifinstr{#1}[
	{%
		\charge_f#1=%
	}
	{%
		\charge_f#1[]=%
	}%
}
\def\charge_f#1[#2]={%
	\CF_ifinstr{#1}:
	{%
		\charge_g#1[#2]=%
	}
	{%
		\charge_g#1:0pt[#2]=%
	}%
}
\def\charge_g#1:#2[#3]=#4,{% #1=angle, #2=offset, #3=code tikz charge, #4=charge
	\CF_stripsp{#1}\CF_ifinteger
	{%
		\pgfmathsetmacro\chargeangle{mod(#1,360)}%
	}
	{%
		\pgfmathanglebetweenpoints{\pgfpointanchor{atom}{center}}{\CF_stripsp{#1}{\pgfpointanchor{atom}}}%
		\let\chargeangle\pgfmathresult% incorrect si (atom.center==atom.#1) && (extra sep==0) TODO: mettre un warning ?
	}%
	\edef\CF_offset{\the\dimexpr#2+0pt}%
	\CF_stripsp{#1}{\CF_distancebetweenpoints{atom}{center}{atom}}\CF_chargedistance
	\CF_eeafter{\node[anchor=center,}{\ifboolKV[charge]{overlay}{overlay,}{}}#3]%
		at([shift=(\chargeangle:\CF_chargedistance pt+\CF_offset)]atom.center){\useKV[charge]{macro charge}{#4}};%
	\charge_d
}

\def\Chembelow{%
	\begingroup
		\let\CF_temp\CF_gobarg
		\CF_chembelowa
}

\def\chembelow{%
	\begingroup
		\let\CF_temp\CF_id
		\CF_chembelowa
}

\def\CF_chembelowa{\CF_testopt\CF_chembelowb\CF_stacksep}

\def\CF_chembelowb[#1]#2#3{%
		\setbox\CF_box\hbox{\printatom{#2}}%
		\expandafter\vtop\CF_temp{to\ht\CF_box}{%
			\offinterlineskip
			\hbox{\printatom{#2}}%
			\kern#1\relax
			\hbox to\wd\CF_box{\hss\printatom{#3}\hss}%
			\CF_temp\vss
			}%
	\endgroup
}

\def\Chemabove{%
	\begingroup
		\let\CF_temp\CF_gobarg
		\CF_chemabovea
}

\def\chemabove{%
	\begingroup
		\let\CF_temp\CF_id
		\CF_chemabovea
}

\def\CF_chemabovea{\CF_testopt\CF_chemaboveb\CF_stacksep}

\def\CF_chemaboveb[#1]#2#3{%
		\setbox\CF_box\hbox{\printatom{#2}}%
		\expandafter\vbox\CF_temp{to\ht\CF_box}{%
			\offinterlineskip
			\CF_temp\vss
			\hbox to\wd\CF_box{\hss\printatom{#3}\hss}%
			\kern#1\relax
			\hbox{\printatom{#2}}%
		}%
	\endgroup
}

\def\chemmove{\CF_testopt\CF_chemmove{}}

\def\CF_chemmove[#1]#2{%
	\CF_doifnotempty{#2}%
	{%
		\expandafter\tikzpicture\expanded{[overlay,remember picture,-CF\CF_ifempty{#1}{}{,\unexpanded{#1}}]}%
			#2%
		\endtikzpicture
	}%
}

\def\chemnameinit#1{%
	\setbox\CF_box_name\hbox{#1}%
	\ifboolKV[chemfig]{gchemname}
		\xdef
		\edef
		     \CF_dpmax{\the\dp\CF_box_name}%
}
\chemnameinit{}

\def\CF_parsemolname#1\\#2\_nil{%
	\hbox to\CF_wdstuffbox{\hss#1\hss}%
	\CF_doifnotempty{#2}{\CF_parsemolname#2\_nil}%
}

\def\chemname{%
	\CF_ifstar
	{%
		\CF_adjust_name_dpfalse
		\CF_chemname_a
	}
	{%
		\CF_adjust_name_dptrue
		\CF_chemname_a
	}%
}

\def\CF_chemname_a{\CF_testopt\CF_chemname_b{}}

\def\CF_chemname_b[#1]#2#3{%
	\CF_stripsp{#1}\CF_ifempty% v1.71 : mise en cohérence avec \name[dimension]{nom}
	{%
		\def\CF_current_name_sep{\useKV[chemfig]{name sep}}%
	}
	{%
		\def\CF_current_name_sep{#1}% 
	}%
	\setbox\CF_box_name\hbox{#2}%
	\edef\CF_wdstuffbox{\the\wd\CF_box_name}%
	\edef\CF_dpstuffbox{\the\dp\CF_box_name}%
	\leavevmode
	\ifdim\CF_dpmax<\CF_dpstuffbox
		\ifboolKV[chemfig]{gchemname}\global{}\let\CF_dpmax\CF_dpstuffbox
	\fi
	\vtop{%
		\box\CF_box_name
		\nointerlineskip
		\kern\dimexpr\CF_current_name_sep\ifCF_adjust_name_dp+\CF_dpmax-\CF_dpstuffbox\fi\relax
		\CF_parsemolname#3\\\_nil
	}%
}
%%%%%%%%%%%%%%%%%%%%%%%%%%%%%%%%%%%%%%%%%%%%%%%%%%%%%%%%%%%%%%%%%%%%%%
%%%        R É A C T I O N S       H O R I Z O N T A L E S         %%%
%%%%%%%%%%%%%%%%%%%%%%%%%%%%%%%%%%%%%%%%%%%%%%%%%%%%%%%%%%%%%%%%%%%%%%

\def\CF_harrowdisplaylabel#1#2#3#4#5#6#7#8{}
\def\hreac{%
	\begingroup
	\CF_testopt\CF_hreac_a{}%
}
\def\CF_hreac_a[#1]{% #1=arg optionnel
	\CF_doifnotempty{#1}{\setKV[chemfig]{#1}}%
	\edef\CF_harrow_length{\the\dimexpr\CF_harrow_min_width*\CF_arrowlength}%
	\CF_cnt_hreac0 % compteur de composé
	\pgfmathsetmacro\CF_arrowdoublesep{\CF_arrowdoublesep/2}%
	\pgfmathsetmacro\CF_arrowdoubleposstart{(1-\CF_arrowdoubleposstart)/2}%
	\pgfmathsetmacro\CF_arrowdoubleposend{1-\CF_arrowdoubleposstart}%
	\let\CF_arrowshiftnodes\CF_gobarg %              neutralisation des macros
	\let\CF_arrowdisplaylabel\CF_harrowdisplaylabel% inutiles de \CF_arrow
	\let\CF_hreac_name_list\empty
	\let\CF_hreac_current_name\empty
	\def\CF_hreac_max_depth{0pt}%
	\def\CF_current_shiftup{0pt}%
	\def\CF_current_shiftup_back{0pt}%
	\def\CF_current_shiftleft{0pt}%
	\edef\CF_current_node_name{c\number\CF_cnt_hreac}%
	\let\CF_hreac_current_compound\empty
	\let\chemname\CF_hreac_chemname % reprogrammation de \chemname
	\let\name\CF_hreac_name
	\CF_compound_is_chemfigfalse
	\ifboolKV[chemfig]{hreac debug}
	{%
		\tikzpicture[every node/.style={draw,anchor=base,inner sep=0pt,outer sep=0pt,minimum size=1pt},baseline]%
	}
	{%
		\tikzpicture[every node/.style={anchor=base,inner sep=0pt,outer sep=0pt,minimum size=0pt},baseline]%
	}%
	\node[anchor=\CF_hreac_init_anchor](\CF_current_node_name){\vphantom{A}};%
	\CF_extracty{\CF_current_node_name}{base}\CF_hreac_posy_initnode
	\CF_hreac_b
}
\def\CF_hreac_b{%
	\futurelet\CF_toksa\CF_hreac_c
}
\def\CF_hreac_c{%
	\CF_ifxcase\CF_toksa
		\endhreac   {}% exécuter \endhreac
		\bgroup     {%
					\ifCF_compound_is_chemfig
						\edef\CF_restore_hashcatcode{\catcode\number`\# =\number\catcode`\# \relax}%
						\catcode`\# 12
					\fi
					\CF_hreac_add_arg
					}%
		\CF_sptoken \CF_hreac_gob_space
		+           \CF_hreac_next_plus
		>           \CF_hreac_next_arrow
		\^          \CF_hreac_set_shiftup
		\>          \CF_hreac_set_shiftleft
		\chemfig    \CF_hreac_next_chemfig
		\chemname   {}% exécuter \CF_hreac_name
		\name       {}% exécuter \CF_hreac_name
		\hreac      \CF_hreac_illegal_command
		\schemestart\CF_hreac_illegal_command
	\CF_ifxcase_elseif
		\CF_hreac_addtok
	\CF_ifxcase_endif
}
\def\CF_hreac_illegal_command#1{%
	\CF_error{Macro '\string#1' not allowed in \string\hreac...\string\endhreac}%
}
\def\CF_hreac_set_shiftup\^{%
	\CF_ifstar
	{%
		\CF_hreac_set_shiftup_a1%
	}
	{%
		\CF_hreac_set_shiftup_a0%
	}%
}
\def\CF_hreac_set_shiftup_a#1#2{% #1=1 ou 0 selon *   #2=dimension
	\edef\CF_current_shiftup{\CF_current_shiftup+#2}%
	\edef\CF_current_shiftup_back{\ifnum#1=1 0pt\else-#2\fi}%
	\CF_hreac_b
}
\def\CF_hreac_set_shiftleft\>#1{%
	\def\CF_current_shiftleft{#1}%
	\CF_hreac_b
}
\def\CF_hreac_addtok#1{%
	\CF_addtomacro\CF_hreac_current_compound{#1}%
	\CF_hreac_b
}
\def\CF_hreac_add_arg#1{%
	\CF_addtomacro\CF_hreac_current_compound{{#1}}%
	\ifCF_compound_is_chemfig
		\CF_restore_hashcatcode
		\CF_compound_is_chemfigfalse
	\fi
	\CF_hreac_b
}
\def\endhreac{%
	\CF_earg\CF_hreac_compose\CF_hreac_current_compound
	\expandafter\CF_hreac_draw_names\CF_hreac_name_list[0={}{}{}]% placer les noms des composés
	\ifboolKV[chemfig]{hreac debug}
	{%
		\CF_extractx{\CF_current_node_name}{0}\CF_current_width
		\pgfextractx\CF_dim_tmp{\pgfpointanchor{c1}{180}}%
		\edef\CF_current_width{\the\dimexpr\CF_current_width-\CF_dim_tmp\relax}% largeur de la réaction
		\draw[dotted,red,line width=0.5pt](c1.west)--([xshift=\CF_current_width]c1.west);%
	}
	{}%
	\endtikzpicture
	\endgroup
}
\def\CF_hreac_draw_names[#1=#2#3#4]{%
	\CF_ifnum{#1=0 }
	{}
	{%
		\node[anchor=north,at=(c#1.south),yshift=-(\CF_hreac_max_depth-#3+\CF_hreac_name_sep)](name#1)%
			{\hbox to#2{\hss\vbox{\CF_compose_harrow_label#4\\\_nil}\hss}};%
		\ifboolKV[chemfig]{hreac debug}
		{%
			\node[anchor=north,at=(name#1.south),fill=purple!60,overlay,opacity=0.5]{\scriptsize\bfseries name#1};
		}%
		{}%
		\CF_hreac_draw_names
	}%
}
\expandafter\def\expandafter\CF_hreac_gob_space\space{%
	\CF_hreac_b
}
\def\CF_hreac_next_plus#1{%
	\CF_earg\CF_hreac_compose\CF_hreac_current_compound
	\CF_hreac_compose{#1}% #1= '+'
	\CF_hreac_b
}
\def\CF_hreac_next_arrow>{%
	\CF_earg\CF_hreac_compose\CF_hreac_current_compound
	\CF_harrow
}

\def\CF_hreac_next_chemfig#1{% #1 est \chemfig
	\CF_testopt\CF_hreac_next_chemfig_a{}% bugfix v1.71
}
\def\CF_hreac_next_chemfig_a[#1]{%
	\CF_addtomacro\CF_hreac_current_compound{\chemfig[#1]}%
	\CF_compound_is_chemfigtrue
	\CF_hreac_b
}
\def\CF_hreac_chemname{%
	\CF_testopt\CF_hreac_chemname_a{}% fixme v1.71
}
\def\CF_hreac_chemname_a[#1]{%
	\CF_stripsp{#1}\CF_ifempty% v1.71
	{%
		\edef\CF_hreac_chemname_current_sep{\useKV[chemfig]{name sep}}%
	}
	{%
		\def\CF_hreac_chemname_current_sep{#1}% 
	}%
	\begingroup
		\catcode`\# 12
		\CF_hreac_chemname_b
}
\def\CF_hreac_chemname_b#1{% #1=composé lu avec \catcode`\#=12
	\endgroup
	\CF_hreac_chemname_c{#1}%
}
\def\CF_hreac_chemname_c#1#2{% #1=composé  #2=nom
	\endgroup
	\def\CF_hreac_current_name{#2}%
	\CF_hreac_b#1%
}
\def\CF_hreac_name{% si \name{nom} est rencontré
	\CF_testopt\CF_hreac_name_a{}% fixme v1.71
}
\def\CF_hreac_name_a[#1]#2{%
	\CF_stripsp{#1}\CF_ifempty% v1.71
	{%
		\edef\CF_hreac_chemname_current_sep{\useKV[chemfig]{name sep}}%
	}
	{%
		\def\CF_hreac_chemname_current_sep{#1}% 
	}%
	\def\CF_hreac_current_name{#2}%
	\CF_hreac_b
}
\def\CF_harrow{%
	\CF_ifnexttok\bgroup
		{\CF_harrow_a}
		{\CF_harrow_b{->}}%
}
\def\CF_harrow_a#1{%
	\CF_earg\CF_stripsp{\detokenize{#1}}\CF_harrow_b
}
\def\CF_harrow_b#1{%
	\ifcsname CF_arrow(#1)\endcsname% bugfix 1.8
		\def\CF_harrow_type{#1}%
	\else
		\CF_error{Undefined arrow  "#1", "->" used instead}%
		\def\CF_harrow_type{->}%
	\fi
	\CF_harrow_c
}
\def\CF_harrow_c{%
	\CF_testopt\CF_harrow_d{}%
}
\def\CF_harrow_d[#1]{%
	\CF_testopt{\CF_harrow_e[#1]}{}%
}
\def\CF_harrow_e[#1][#2]{%
	\CF_doifnotempty{#1}
	{%
		\setbox\CF_box_harrowup\hbox
		{%
			\vbox{\pgfinterruptpicture\CF_compose_harrow_label#1\\\_nil\endpgfinterruptpicture}%
		}%
		\ifdim\dimexpr\wd\CF_box_harrowup+\CF_harrow_label_xsep*2\relax>\CF_harrow_length\relax
			\edef\CF_harrow_length{\the\dimexpr\wd\CF_box_harrowup+\CF_harrow_label_xsep*2\relax}%
		\fi
	}%
	\CF_doifnotempty{#2}
	{%
		\setbox\CF_box_harrow_down\hbox
		{%
			\vtop{\pgfinterruptpicture\CF_compose_harrow_label#2\\\_nil\endpgfinterruptpicture}%
		}%
		\ifdim\dimexpr\wd\CF_box_harrow_down+\CF_harrow_label_xsep*2\relax>\CF_harrow_length\relax
			\edef\CF_harrow_length{\the\dimexpr\wd\CF_box_harrow_down+\CF_harrow_label_xsep*2\relax}%
		\fi
	}%
	\CF_hreac_compose{\kern\CF_harrow_length}%
	\node[at=(\CF_current_node_name.180)](harrowstart){};%
	\node[at=(\CF_current_node_name.0)](harrowend){};%
	\CF_eesecond{\def\CF_arrowtip}{\expandafter\CF_gobarg\CF_arrowhead}%
	\edef\CF_arrowstartnode{harrowstart}%
	\edef\CF_arrowendnode  {harrowend}%
	\edef\CF_arrowcurrentstyle{% bugfix 1.8
		\unexpanded\expandafter{\CF_arrowhead}%
		\unless\ifx\CF_defaultarrowstyle\empty
			,\unexpanded\expandafter{\CF_defaultarrowstyle}%
		\fi
	}%
	\csname CF_arrow(\CF_harrow_type)\endcsname[][][][][][][][]\_nil
	\CF_doifnotempty{#1}
	{%
		\node[anchor=south,
		      at=(\CF_current_node_name.north),
		      yshift=\CF_arrowlabelsep]%
		      (labelup\the\CF_cnt_hreac)%
		      {\hbox to\dimexpr\CF_harrow_length-\CF_harrow_label_xsep*2\relax
		      	{\CF_harrow_label_align_left\box\CF_box_harrowup\CF_harrow_label_align_right}%
		      };%
		\ifboolKV[chemfig]{hreac debug}
		{%
		\node[anchor=south,
		      at=(labelup\the\CF_cnt_hreac.north),
		      fill=blue!60,
		      overlay,
		      opacity=0.5]%
		      {\scriptsize\bfseries l-sup\the\CF_cnt_hreac};
		}
		{}%
	}%
	\CF_doifnotempty{#2}
	{%
		\node[%
		        anchor=north,
		        at=(\CF_current_node_name.south),
		        yshift=-\CF_arrowlabelsep%
		      ]%
		      (labeldown\the\CF_cnt_hreac)%
		      {%
		      \hbox to\dimexpr\CF_harrow_length-\CF_harrow_label_xsep*2\relax
		      	{\CF_harrow_label_align_left\box\CF_box_harrow_down\CF_harrow_label_align_right}%
		      };%
		\ifboolKV[chemfig]{hreac debug}
		{%
			\node[%
			        anchor=north,
			        at=(labeldown\the\CF_cnt_hreac.south),
			        fill=blue!60,
			        overlay,
			        opacity=0.5
			      ]
			      {\scriptsize\bfseries l-inf\the\CF_cnt_hreac};
		}
		{}%
	}%
	\CF_hreac_b
}
\def\CF_hreac_compose#1{%
	\advance\CF_cnt_hreac1
	\let\CF_previousnodename\CF_current_node_name
	\edef\CF_current_node_name{c\number\CF_cnt_hreac}%
	\edef\CF_hreac_current_compound_sep{\CF_ifnum{\CF_cnt_hreac=1 }{0pt}{\CF_hreac_sep}}%
	\node[%
	        anchor=180,
	        at=(\CF_previousnodename.0),
	        xshift=\CF_hreac_current_compound_sep+\CF_current_shiftleft,
	        yshift=\CF_current_shiftup%
	      ]%
	      (\CF_current_node_name){#1};%
	\ifboolKV[chemfig]{hreac debug}
	{%
		\node[%
		        anchor=south,
		        at=(\CF_current_node_name.north),
		        fill=green!60,
		        overlay,
		        opacity=0.5%
		      ]%
		      {\scriptsize\bfseries \CF_current_node_name};%
		\node[at=(\CF_current_node_name.west),fill,red]{};
	}
	{}%
	\pgfextracty\CF_dim_tmp{\pgfpointanchor{\CF_current_node_name}{south}}%
	\edef\CF_current_depth{\the\dimexpr\CF_hreac_posy_initnode-\CF_dim_tmp\relax}%
	\CF_ifnum{0\ifx\CF_hreac_current_name\empty1\else\ifx\CF_hreac_chemname_current_sep\empty1\fi\fi>0 }
	{% si composé sans nom OU SINON [décalage] non précisé dans \chemname ou \name
		\ifdim\CF_current_depth>\CF_hreac_max_depth\relax
			\let\CF_hreac_max_depth\CF_current_depth
		\fi
	}
	{%
		\ifdim\dimexpr\CF_current_depth+\CF_hreac_chemname_current_sep\relax>\CF_hreac_max_depth\relax
			\edef\CF_hreac_max_depth{\the\dimexpr\CF_current_depth+\CF_hreac_chemname_current_sep-\CF_hreac_name_sep\relax}%
		\fi
	}%
	\CF_ifx\CF_hreac_current_name\empty
	{}
	{%
		\CF_extractx{\CF_current_node_name}{0}\CF_current_width
		\pgfextractx\CF_dim_tmp{\pgfpointanchor{\CF_current_node_name}{180}}%
		\edef\CF_current_width{\the\dimexpr\CF_current_width-\CF_dim_tmp\relax}% largeur du composé
		\CF_xaddtomacro\CF_hreac_name_list{%
			[%
				\the\CF_cnt_hreac=%
				{\CF_current_width}%
				{\CF_current_depth}%
				{\unexpanded\expandafter{\CF_hreac_current_name}}%
			]%
		}%
		\let\CF_hreac_current_name\empty
	}%
	\def\CF_hreac_current_compound{}%
	\def\CF_current_shiftleft{0pt}%
	\let\CF_current_shiftup\CF_current_shiftup_back
	\def\CF_current_shiftup_back{0pt}%
}
\def\CF_compose_harrow_label#1\\#2\_nil{%
	\hbox{#1}%
	\CF_doifnotempty{#2}{\CF_compose_harrow_label#2\_nil}%
}

%%%%%%%%%%%%%%%%%%%%%%%%%%%%%%%%%%%%%%%%%%%%%%%%%%%%%%%%%%%%%%%%%%%%%%
%%%           S C H É M A S      R É A C T I O N N E L S           %%%
%%%%%%%%%%%%%%%%%%%%%%%%%%%%%%%%%%%%%%%%%%%%%%%%%%%%%%%%%%%%%%%%%%%%%%
\let\CF_schemenest\CF_zero

\def\CF_subscheme{\CF_testopt\CF_subscheme_a{}}
\def\CF_subscheme_a[#1]#2{\schemestart[#1]#2\schemestop}

\def\chemleft#1#2\chemright#3{%
	\leavevmode
	\begingroup
		\setbox0\hbox{$\vcenter{\hbox{}}$}\edef\CF_delimmathht{\the\ht0}%
		\setbox0\hbox{#2}\edef\CF_delimdim{\the\dimexpr(\ht0+\dp0)/2}%
		\edef\CF_delimshift{\the\dimexpr(\ht0-\dp0)/2-\CF_delimmathht}%
		\raise\CF_delimshift\hbox{$\left#1\vrule height\CF_delimdim depth\CF_delimdim width0pt\right.$}\box0
		\raise\CF_delimshift\hbox{$\left.\vrule height\CF_delimdim depth\CF_delimdim width0pt\right#3$}%
	\endgroup
}

\def\chemright#1{%
	\CF_warning{"\string\chemright\string#1" ignored! No \string\chemleft\space previously found.}%
}

\def\chemup#1#2\chemdown#3{%
	\begingroup
		\setbox0\hbox{\printatom{#2}}\edef\CF_delimdim{\the\dimexpr\wd0/2}%
		\tikzpicture[every node/.style={inner sep=0pt,outer sep=0pt,minimum size=0pt},baseline]%
			\node[anchor=base west](chem@stuff){\box0};%
			\node[at=(chem@stuff.north),anchor=east,rotate=-90]{$\left#1\vrule height\CF_delimdim depth\CF_delimdim width0pt\right.$};%
			\node[at=(chem@stuff.south),anchor=west,rotate=-90]{$\left.\vrule height\CF_delimdim depth\CF_delimdim width0pt\right#3$};%
		\endtikzpicture
	\endgroup
}

\def\chemdown#1{%
	\CF_warning{"\string\chemdown\string#1" ignored! No \string\chemup\space previously found.}%
}

\def\CF_setstyle#1,#2,#3\_nil#4#5#6{%
	\def#4{#1}%
	\let#5\empty
	\let#6\empty
	\CF_iffirsttokis{#2\relax}\CF_quark
	{}%
	{%
		\def#5{#2}%
		\CF_iffirsttokis{#3\relax}\CF_quark
		{}%
		{%
			\CF_setstyle_a#3\_nil#6%
		}%
	}%
}
\def\CF_setstyle_a#1,\CF_quark#2\_nil#3{\def#3{#1}}

\def\CF_and{\futurelet\CF_toksa\CF_and_a}

\def\CF_and_a{%
	\CF_ifx\CF_toksa\bgroup
	{%
		\CF_and_b
	}
	{%
		\CF_and_b{}%
	}%
}

\def\CF_and_b#1{%
	\CF_setstyle#1,\CF_quark,\CF_quark\_nil\CF_signspaceante_\CF_signspacepost_\CF_signvshift_
	\CF_doifnotempty\CF_signspaceante_
	{%
		\let\CF_signspaceante\CF_signspaceante_
	}%
	\CF_doifnotempty\CF_signspacepost_
	{%
		\let\CF_signspacepost\CF_signspacepost_
	}%
	\CF_doifnotempty\CF_signvshift_
	{%
		\let\CF_signvshift\CF_signvshift_
	}%
	\raise\CF_signvshift\hbox{\kern\CF_signspaceante$+$\kern\CF_signspacepost}%
}

\def\schemestart{%
	\begingroup
	\CF_ifnum{\CF_schemenest=0 }
	{%
		\useKV[chemfig]{schemestart code}%
	}%
	{}%
	\xdef\CF_schemenest{\number\numexpr\CF_schemenest+1}%
	\CF_testopt\CF_schemestart_a{}%
}

\def\CF_schemestart_a[#1]{%
	\setchemfig{#1}%
	\CF_eesecond{\def\CF_arrowtip}{\expandafter\CF_gobarg\CF_arrowhead}%
	\edef\CF_defaultarrowstyle{\unexpanded\expandafter{\CF_arrowhead,}\unexpanded\expandafter{\CF_defaultarrowstyle}}%
	\let\CF_arrowstyle\CF_defaultarrowstyle
	\pgfmathsetmacro\CF_arrowdoublesep{\CF_arrowdoublesep/2}%
	\pgfmathsetmacro\CF_arrowdoubleposstart{(1-\CF_arrowdoubleposstart)/2}%
	\pgfmathsetmacro\CF_arrowdoubleposend{1-\CF_arrowdoubleposstart}%
	\ifboolKV[chemfig]{scheme debug}
	{%
		\tikzpicture[every node/.style={draw,anchor=base,inner sep=0pt,outer sep=0pt,minimum size=1.5pt},baseline,remember picture]%
	}
	{%
		\tikzpicture[every node/.style={anchor=base,inner sep=0pt,outer sep=0pt,minimum size=0pt},baseline,remember picture]%
	}%
	\let\merge\CF_merge
	\expandafter\let\csname+\endcsname\CF_and
	\let\arrow\CF_arrow
	\let\schemestop\CF_schemestop
	\let\subscheme\CF_subscheme
	\ifnum\CF_schemenest=1 % la commande n'est pas imbriquée ?
		\CF_cnt_compound0
	\fi
	\edef\CF_currentnodename{c\number\CF_cnt_compound}%
	\let\CF_nextnodename\empty
	\let\CF_nextnodestyle\empty
	\let\CF_directarrowlist\empty
	\ifboolKV[chemfig]{scheme debug}
	{%
		\node[fill,green](\CF_currentnodename){};%
	}
	{%
		\node(\CF_currentnodename){};%
	}%
	\let\CF_compound\empty
	\CF_schemestart_c
}

\def\CF_schemestart_c{%
	\futurelet\CF_toksa\CF_schemestart_e
}

\expandafter\def\expandafter\CF_schemestart_d\space{\futurelet\CF_toksa\CF_schemestart_e}

\def\CF_schemestart_e{% ... et l'examine :
	\CF_iffirsttokin_a{\arrow\schemestop\merge}%
	{}
	{%
		\CF_ifx\CF_toksa\bgroup
		{%
			\ifCF_compound_is_chemfig% bugfix 1.6
				\edef\CF_restore_hashcatcode{\catcode\number`\#=\number\catcode`\# \relax}%
				\catcode`\#12
			\fi
			\CF_addnextarg
		}
		{%
			\CF_ifx\CF_toksa\CF_sptoken
			{%
				\CF_addtomacro\CF_compound{ }%
				\CF_schemestart_d
			}
			{%
				\CF_ifx\CF_toksa\chemfig
					\CF_compound_is_chemfigtrue% mettre le flag à vrai
					{}%
				\afterassignment\CF_schemestart_c
				\CF_addtomacro\CF_compound
			}%
		}%
	}%
}

\def\CF_addnextarg#1{%
	\CF_addtomacro\CF_compound{{#1}}%
	\ifCF_compound_is_chemfig% bugfix 1.6
		\CF_restore_hashcatcode
		\CF_compound_is_chemfigfalse% mettre le flag à faux
	\fi
	\CF_schemestart_c
}

\def\CF_displaycompound#1#2{% #1 = nom et #2 = style
	\CF_doifnotempty\CF_compound
	{%
		\global\advance\CF_cnt_compound1
		\CF_ifx\CF_defaultcompoundstyle\empty
		{%
			\let\CF_currentnodestyle\empty
		}
		{%
			\expandafter\def\expandafter\CF_currentnodestyle\expandafter{\CF_defaultcompoundstyle,}%
		}%
		\CF_addtomacro\CF_currentnodestyle{anchor=\CF_nextnodeanchor,at=(\CF_currentnodename)}%
		\CF_ifempty{#2}%
		{%
			\CF_doifnotempty\CF_nextnodestyle
			{%
				\CF_eaddtomacro\CF_currentnodestyle{\expandafter,\CF_nextnodestyle}%
			}%
		}
		{%
			\CF_doifnotempty\CF_nextnodestyle
			{%
				\CF_warning{two styles for the same node, first style "\CF_nextnodestyle" ignored}%
			}%
			\CF_addtomacro\CF_currentnodestyle{,#2}%
		}%
		\CF_ifempty{#1}
		{%
			\edef\CF_temp{%
				\CF_ifempty\CF_nextnodename
				{%
					c\number\CF_cnt_compound
				}
				{%
					\CF_nextnodename
				}%
			}%
		}
		{%
			\CF_doifnotempty\CF_nextnodename
			{%
				\CF_warning{two names for the same node, first name "\CF_nextnodename" ignored}%
			}%
			\edef\CF_temp{#1}%
		}%
		\CF_eafter{\node[}\CF_currentnodestyle](\CF_temp){\CF_compound};%
		\ifboolKV[chemfig]{scheme debug}%
		{%
			\node[draw=none,anchor=270,at=(\CF_temp.90),fill=green!60,overlay,opacity=0.5]{\scriptsize\bfseries\CF_temp};%
		}
		{}%
		\let\CF_currentnodename\CF_temp
	}%
}

\def\CF_schemestop{%
	\CF_displaycompound{}{}%
	\CF_directarrowlist
	\endtikzpicture
	\xdef\CF_schemenest{\number\numexpr\CF_schemenest-1}%
	\CF_ifnum{\CF_schemenest=0 }
	{%
		\useKV[chemfig]{schemestop code}%
	}%
	{}%
	\endgroup
}

\def\CF_analysearrowarg#1{\CF_analysearrowarg_a#1[]\_nil}

\def\CF_analysearrowarg_a#1[#2]#3\_nil{%
	\CF_ifinstr{#1}.
	{%
		\CF_addtomacro\CF_temp{#1[#2]}%
	}
	{%
		\CF_addtomacro\CF_temp{#1.[#2]}%
	}%
}

\def\CF_arrow{%
	\CF_ifnexttok(%
	{%
		\CF_arrowa
	}
	{%
		\CF_ifnexttok\bgroup
		{%
			\CF_arrowb(.[]--.[])%
		}
		{%
			\CF_arrowb(.[]--.[]){}%
		}%
	}%
}

\def\CF_arrowa(#1--#2){%
	\def\CF_temp{(}%
	\CF_analysearrowarg{#1}%
	\CF_addtomacro\CF_temp{--}%
	\CF_analysearrowarg{#2}%
	\CF_addtomacro\CF_temp)%
	\CF_ifnexttok\bgroup
	{%
		\expandafter\CF_arrowb\CF_temp
	}
	{%
		\expandafter\CF_arrowb\CF_temp{}%
	}%
}

\def\CF_arrowb(#1.#2[#3]--#4.#5[#6])#7{%
	\def\CF_currentarrowtype{#7}% nom de la flèche
	\CF_doifempty\CF_currentarrowtype
	{%
		\def\CF_currentarrowtype{->}%
	}%
	\edef\CF_currentarrowname{\expandafter\CF_extract_arrowname\CF_currentarrowtype[\_nil}%
	\CF_testopt{\CF_arrowc(#1.#2[#3]--#4.#5[#6])}{}%
}

\def\CF_arrowc(#1.#2[#3]--#4.#5[#6])[#7]{%
	\def\CF_temp{\CF_arrowe(#1.#2[#3]--#4.#5[#6])}%
	\CF_arrowd#7,\empty,\empty\_nil
}

\def\CF_arrowd#1,#2,#3\_nil{%
	\CF_addtomacro\CF_temp{{#1}}%
	\CF_eaddtomacro\CF_temp{\expandafter{#2}}%
	\expandafter\CF_eaddtomacro\expandafter\CF_temp\expandafter{\expandafter\expandafter\expandafter{\expandafter\CF_sanitizelastitem#3,\empty\_nil}}%
	\CF_temp
}

% #1, #4 : nom des nodes    #2, #5 : ancres des nodes    #3, #6 : styles des nodes
% #7 : angle flèche    #8 : longueur flèche    #9 : style tikz de la flèche
\def\CF_arrowe(#1.#2[#3]--#4.#5[#6])#7#8#9{%
	\let\CF_arrowcurrentstyle\CF_arrowstyle
	\if @\expandafter\CF_firsttonil\detokenize{#1.}\_nil% si #1 commence par @
		\if @\expandafter\CF_firsttonil\detokenize{#4.}\_nil
			\CF_eaddtomacro\CF_directarrowlist{\expandafter\CF_directarrow\expandafter{\CF_currentarrowtype}{#1}{#2}{#4}{#5}{#9}}%
			\let\CF_nextaction\CF_schemestart_c
		\else
			\CF_doifnotempty\CF_arrowcurrentstyle
			{%
				\CF_addtomacro\CF_arrowcurrentstyle,%
			}%
			\CF_doifnotempty{#9}
			{%
				\CF_addtomacro\CF_arrowcurrentstyle{#9,}%
			}%
			\CF_doifnotempty{#3}% bugfix 1.8
			{%
				\CF_warning{Node "#1" already drown, style "\detokenize{#3}" ignored}%
			}%
			\CF_esecond{\CF_displaycompound{}}\CF_nextnodestyle% bugfix 1.8
			\def\CF_nextnodename{#4}%
			\expandafter\def\expandafter\CF_currentnodename\expandafter{\CF_gobarg#1}%
			\let\CF_arrowstartname\CF_currentnodename
			\let\CF_arrowendname\CF_nextnodename
			\CF_arrow_f{#7}{#8}{#2}{#5}%
			\def\CF_nextnodestyle{#6}%
		\fi
	\else
		\CF_doifnotempty\CF_arrowcurrentstyle
		{%
			\CF_addtomacro\CF_arrowcurrentstyle,%
		}%
		\CF_doifnotempty{#9}
		{%
			\CF_addtomacro\CF_arrowcurrentstyle{#9,}%
		}%
		\if @\expandafter\CF_firsttonil\detokenize{#2.}\_nil
			\CF_error{syntax "(<name>--@<name>)" is not allowed}%
		\else
			\CF_displaycompound{#1}{#3}%
			\edef\CF_arrowstartname{%
				\CF_ifempty{#1}
					\CF_currentnodename
					{#1}%
				\CF_doifnotempty{#2}{.#2}%
			}%
			\CF_arrow_f{#7}{#8}{#2}{#5}%
			\def\CF_nextnodename{#4}%
			\def\CF_nextnodestyle{#6}%
		\fi
	\fi
	\CF_arrowgobspaces% mange les espaces puis exécute \CF_nextaction
}

\def\CF_arrowgobspaces{\futurelet\CF_toksa\CF_arrowgobspaces_a}

\def\CF_arrowgobspaces_a{%
	\CF_ifx\CF_sptoken\CF_toksa
		\CF_arrowgobspaces_b
		\CF_nextaction
}

\expandafter\def\expandafter\CF_arrowgobspaces_b\space{\futurelet\CF_toksa\CF_arrowgobspaces_a}

\def\CF_arrow_f#1#2#3#4{% #1=angle   #2=longueur    #3=ancre départ   #4=ancre arrivée
	\def\CF_nextaction{\let\CF_compound\empty\CF_schemestart_c}%
	\def\CF_arrowcurrentangle{#1}\CF_doifempty\CF_arrowcurrentangle{\let\CF_arrowcurrentangle\CF_arrowangle}%
	\def\CF_currentarrowlength{#2}\CF_doifempty\CF_currentarrowlength{\let\CF_currentarrowlength\CF_arrowlength}%
	\node[at=(\CF_currentnodename.\CF_ifempty{#3}\CF_arrowcurrentangle{#3}),shift=(\CF_arrowcurrentangle:\CF_currentarrowlength*\CF_compoundsep),cyan,fill](end@arrow@i@\number\CF_schemenest){};%
	\edef\CF_arrowendname{end@arrow@i@\number\CF_schemenest\CF_doifnotempty{#4}{.#4}}%
	\ifboolKV[chemfig]{scheme debug}
	{%
		\node[at=(\CF_currentnodename.\CF_ifempty{#3}\CF_arrowcurrentangle{#3}),shift=(\CF_arrowcurrentangle:\CF_arrowoffset),red,fill](start@arrow){};%
		\node[at=(\CF_currentnodename.\CF_ifempty{#3}\CF_arrowcurrentangle{#3}),shift=(\CF_arrowcurrentangle:\CF_currentarrowlength*\CF_compoundsep-\CF_arrowoffset),red,fill](end@arrow){};%
	}
	{%
		\node[at=(\CF_currentnodename.\CF_ifempty{#3}\CF_arrowcurrentangle{#3}),shift=(\CF_arrowcurrentangle:\CF_arrowoffset)](start@arrow){};%
		\node[at=(\CF_currentnodename.\CF_ifempty{#3}\CF_arrowcurrentangle{#3}),shift=(\CF_arrowcurrentangle:\CF_currentarrowlength*\CF_compoundsep-\CF_arrowoffset)](end@arrow){};%
	}%
	\def\CF_arrowstartnode{start@arrow}\def\CF_arrowendnode{end@arrow}%
	\unless\ifcsname CF_arrow(\CF_currentarrowname)\endcsname
		\CF_error{Undefined arrow "\detokenize\expandafter{\CF_currentarrowname}", "->" used instead}%
		\def\CF_currentarrowtype{->}%
	\fi
	\csname CF_arrow(\CF_currentarrowname)\expandafter\expandafter\expandafter\endcsname
		\expandafter\CF_grabarrowargs\CF_currentarrowtype[]\_nil[][][][][][][][]\_nil
	\def\CF_currentnodename{end@arrow@i@\number\CF_schemenest}%
	\edef\CF_nextnodeanchor{\CF_ifempty{#4}{180+\CF_arrowcurrentangle}{#4}}%
}

% trace un flèche initiée par (@nom--@nom)
% #1=type de flèche   #2=nom depart   #3=ancre départ    #4=nom arrivée    #5=ancre arrivée    #6=style flèche
\def\CF_directarrow#1#2#3#4#5#6{%
	\edef\CF_arrowstartname{\CF_gobarg#2\ifx\empty#3\empty\else.#3\fi}% BUGFIX 1.6d
	\edef\CF_arrowendname{\CF_gobarg#4\ifx\empty#5\empty\else.#5\fi}% BUGFIX 1.6d
	\path[sloped,allow upside down](\CF_gobarg#2\ifx\empty#3\empty\else.#3\fi)--(\CF_gobarg#4\ifx\empty#5\empty\else.#5\fi)%
		coordinate[pos=0,xshift=\CF_arrowoffset](start@direct@arrow)%
		coordinate[pos=1,xshift=-\CF_arrowoffset](end@direct@arrow);%
	\def\CF_arrowstartnode{start@direct@arrow}%
	\def\CF_arrowendnode{end@direct@arrow}%
	\pgfmathanglebetweenpoints
		{\pgfpointanchor{\CF_gobarg#2}{\ifx\empty#3\empty center\else#3\fi}}% Ne pas utiliser \CF_ifempty ici !!!
		{\pgfpointanchor{\CF_gobarg#4}{\ifx\empty#5\empty center\else#5\fi}}%
	\let\CF_arrowcurrentangle\pgfmathresult
	\CF_doifnotempty{#6}{\CF_addtomacro\CF_arrowcurrentstyle{#6,}}%
	\csname\CF_grabarrowname#1[]\_nil\expandafter\endcsname\CF_grabarrowargs#1[]\_nil[][][][][][][][]\_nil
}

\def\CF_mergegrabchardir#1[#2][#3]#4\_nil{%
	\CF_eafter{\futurelet\CF_toksa\CF_gobtonil}{\CF_firsttonil#1>\_nil}\_nil
	\ifx>\CF_toksa
		\def\CF_merge_angle{0}\def\CF_mergeextreme{xmax}\def\CF_mergesign{+}%
	\else
	\ifx<\CF_toksa
		\def\CF_merge_angle{180}\def\CF_mergeextreme{xmin}\def\CF_mergesign{-}%
	\else
	\ifx^\CF_toksa
		\def\CF_merge_angle{90}\def\CF_mergeextreme{ymax}\def\CF_mergesign{+}%
	\else
	\ifx v\CF_toksa
		\def\CF_merge_angle{-90}\def\CF_mergeextreme{ymin}\def\CF_mergesign{-}%
	\fi\fi\fi\fi
	\def\CF_mergelabelup{#2}\def\CF_mergelabeldo{#3}%
}

\def\CF_merge#1({%
	\CF_mergegrabchardir#1[][]\_nil
	\CF_merge_a(%
}

\def\CF_merge_a#1--(#2){\CF_testopt{\CF_merge_b#1--(#2)}{}}

\def\CF_merge_b#1--(#2)[#3]{%
	\CF_displaycompound{}{}%
	\CF_parsemergeopt#3,\CF_quark,\CF_quark,\CF_quark\_nil
	\def\CF_mergexmax{-16383.99999pt}\let\CF_mergeymax\CF_mergexmax
	\def\CF_mergexmin{16383.99999pt}\let\CF_mergeymin\CF_mergexmin
	\CF_mergeparsenodelist#1(\relax)% calcule les maxi des positions
	\pgfmathsetmacro\CF_mergeextremeresult{\csname CF_merge\CF_mergeextreme\endcsname\CF_mergesign\CF_mergefromcoeff*\CF_compoundsep}%
	\CF_merge_c#1(\relax)% trace les lignes entre les noeuds précédents et la ligne de jonction
	\expandafter\def\expandafter\CF_temp\expandafter{\expandafter[\CF_mergestyle,shorten <=0,shorten >=0,-]}%
	\if x\expandafter\CF_firsttonil\CF_mergeextreme\_nil
		\CF_addtomacro\CF_temp{(\CF_mergeextremeresult pt,\CF_mergeymax)--(\CF_mergeextremeresult pt,\CF_mergeymin)}%
	\else
		\CF_addtomacro\CF_temp{(\CF_mergexmin,\CF_mergeextremeresult pt)--(\CF_mergexmax,\CF_mergeextremeresult pt)}%
	\fi
	\expandafter\draw\CF_temp node[pos=\CF_mergesplitcoeff](merge@point){}% trace la ligne de jonction
		node[at=(merge@point),shift=(\CF_merge_angle:\CF_compoundsep*\CF_mergetocoeff-\CF_arrowoffset)](end@merge){}%
		node[at=(merge@point),shift=(\CF_merge_angle:\CF_compoundsep*\CF_mergetocoeff)](end@merge@i){};%
	\let\CF_arrowcurrentangle\CF_merge_angle
	\CF_eafter{\draw[}\CF_mergestyle,shorten <=0](merge@point)--(end@merge)%
		\expandafter\CF_arrowdisplaylabel_a\expandafter{\CF_mergelabelup}{.5}+\expandafter\CF_arrowdisplaylabel_a\expandafter{\CF_mergelabeldo}{.5}-;%
	\def\CF_currentnodename{end@merge@i}%
	\let\CF_temp\empty
	\CF_analysearrowarg{#2}%
	\expandafter\CF_merge_d\CF_temp\_nil
}

\def\CF_merge_c(#1){%
	\if\relax\expandafter\noexpand\CF_firsttonil#1\_nil
	\else
		\CF_ifdot{#1}%
		{%
			\edef\merge_currentnodename{\CF_beforedot#1\_nil}%
			\edef\merge_currentanchor{\CF_afterdot#1\_nil}%
		}%
		{%
			\def\merge_currentnodename{#1}%
			\let\merge_currentanchor\CF_merge_angle
		}%
		\if x\expandafter\CF_firsttonil\CF_mergeextreme\_nil
			\pgfextracty\CF_dim_tmp{\pgfpointanchor\merge_currentnodename\merge_currentanchor}%
			\CF_eafter{\draw[}\CF_mergestyle,shorten >=0,-]([shift=(\CF_merge_angle:\CF_arrowoffset)]\merge_currentnodename.\merge_currentanchor)--(\CF_mergeextremeresult pt,\CF_dim_tmp);%
		\else
			\pgfextractx\CF_dim_tmp{\pgfpointanchor\merge_currentnodename\merge_currentanchor}%
			\CF_eafter{\draw[}\CF_mergestyle,shorten >=0,-]([shift=(\CF_merge_angle:\CF_arrowoffset)]\merge_currentnodename.\merge_currentanchor)--(\CF_dim_tmp,\CF_mergeextremeresult pt);%
		\fi
		\expandafter\CF_merge_c
	\fi
}

\def\CF_merge_d#1.#2[#3]\_nil{%
	\def\CF_nextnodename{#1}%
	\edef\CF_nextnodeanchor{%
		\CF_ifempty{#2}
		{%
			180+\CF_merge_angle
		}
		{%
			#2%
		}%
	}%
	\def\CF_nextnodestyle{#3}%
	\let\CF_compound\empty
	\CF_schemestart_c
}

\def\CF_parsemergeopt#1,#2,#3,#4\_nil{%
	\CF_ifempty{#1}
	{%
		\def\CF_mergefromcoeff{0.5}%
	}
	{%
		\def\CF_mergefromcoeff{#1}%
	}%
	\def\CF_mergetocoeff{0.5}%
	\def\CF_mergesplitcoeff{0.5}%
	\expandafter\def\expandafter\CF_mergestyle\expandafter{\CF_arrowhead}%
	\CF_iffirsttokis{#2\relax}\CF_quark
	{}
	{%
		\CF_ifempty{#2}
		{%
			\def\CF_mergetocoeff{0.5}%
		}
		{%
			\def\CF_mergetocoeff{#2}%
		}%
		\CF_iffirsttokis{#3\relax}\CF_quark
		{}
		{%
			\CF_ifempty{#3}
			{%
				\def\CF_mergesplitcoeff{0.5}%
			}
			{%
				\def\CF_mergesplitcoeff{#3}%
			}%
			\CF_iffirsttokis{#4\relax}\CF_quark
			{}
			{%
				\CF_parsemergeopt_a#4\_nil
			}%
		}%
	}%
}

\def\CF_parsemergeopt_a#1,\CF_quark#2\_nil{%
	\CF_ifempty{#1}
	{}
	{%
		\CF_addtomacro\CF_mergestyle{,#1}%
	}%
}

\def\CF_mergeparsenodelist(#1){%
	\if\relax\expandafter\noexpand\CF_firsttonil#1\_nil
	\else
		\CF_ifdot{#1}%
		{%
			\edef\merge_currentnodename{\CF_beforedot#1\_nil}%
			\edef\merge_currentanchor{\CF_afterdot#1\_nil}%
		}%
		{%
			\def\merge_currentnodename{#1}%
			\let\merge_currentanchor\CF_merge_angle
		}%
		\pgfextractx\CF_dim_tmp{\pgfpointanchor\merge_currentnodename\merge_currentanchor}%
		\ifdim\CF_dim_tmp>\CF_mergexmax
			\edef\CF_mergexmax{\the\CF_dim_tmp}%
		\fi
		\ifdim\CF_dim_tmp<\CF_mergexmin
			\edef\CF_mergexmin{\the\CF_dim_tmp}%
		\fi
		\pgfextracty\CF_dim_tmp{\pgfpointanchor\merge_currentnodename\merge_currentanchor}%
		\ifdim\CF_dim_tmp>\CF_mergeymax
			\edef\CF_mergeymax{\the\CF_dim_tmp}%
		\fi
		\ifdim\CF_dim_tmp<\CF_mergeymin
			\edef\CF_mergeymin{\the\CF_dim_tmp}%
		\fi
		\expandafter\CF_mergeparsenodelist
	\fi
}

\def\CF_extract_arrowname#1[#2\_nil{\detokenize{#1}}%
\def\CF_grabarrowname#1[#2\_nil{\detokenize{CF_arrow(#1)}}
\def\CF_grabarrowargs#1[#2\_nil{[#2}

\def\CF_makeparametertext#1{%
	\toks0{}%
	\CF_cnt_group#1\relax
	\CF_makeparametertext_a1%
}

\def\CF_makeparametertext_a#1{%
	\unless\ifnum#1>\CF_cnt_group
		\toks0\expandafter{\the\toks0[###1]}%
		\expandafter\CF_makeparametertext_a\expandafter{\number\numexpr#1+1\expandafter}%
	\fi
}

% #1 est le nombre d'arguments optionnels, #2 est le nom et #3 le code
\def\definearrow#1#2#3{%
	\begingroup
		\CF_makeparametertext{#1}%
	\expandafter\endgroup
	\expandafter\def\csname\detokenize{CF_arrow(#2)}\expandafter\endcsname\the\toks0{#3\CF_gobtonil}%
}

\def\CF_ifdot#1{\CF_ifdot_a#1.\_nil}
\def\CF_ifdot_a#1.#2\_nil{\ifx\empty#2\empty\expandafter\CF_execsecond\else\expandafter\CF_execfirst\fi}
\def\CF_beforedot#1.#2\_nil{#1}
\def\CF_afterdot#1.#2\_nil{#2}

\def\CF_rotatenode*#1#2\_nil{%
	\CF_ifdot{#1}
	{%
		\CF_beforedot#1\_nil
	}
	{%
		#1%
	}%
}
\def\CF_anchornode*#1#2\_nil#3{%
	\CF_ifdot{#1}
	{%
		\CF_afterdot#1\_nil
	}
	{%
		\CF_arrowcurrentangle-#390-#1%
	}%
}

% #1 = label  #2 = position  #3 = + ou - (au dessus ou au dessous)  #4 : nom du noeud de départ
% #5 = label  #6 = position  #7 = + ou - (au dessus ou au dessous)  #8 : nom du noeud de fin
\def\CF_arrowdisplaylabel#1#2#3#4#5#6#7#8{%
	\CF_doifnotempty{#1#5}
		{\path(#4)--(#8)\CF_arrowdisplaylabel_a{#1}{#2}{#3}\CF_arrowdisplaylabel_a{#5}{#6}{#7};}%
}

\def\CF_arrowdisplaylabel_a#1#2#3{%
	\CF_doifnotempty{#1}
		{%
		\if*\expandafter\CF_firsttonil\detokenize{#1}\_nil
			\ifboolKV[chemfig]{scheme debug}
			{%
				node[pos=#2,sloped,yshift=#3\CF_arrowlabelsep,draw,fill,cyan](shifted@node){}%
				node[draw,rotate=\CF_rotatenode#1\_nil,anchor=\CF_anchornode#1\_nil#3,at=(shifted@node)]{\expandafter\CF_gobarg\CF_gobarg#1}%
			}
			{%
				node[pos=#2,sloped,yshift=#3\CF_arrowlabelsep](shifted@node){}%
				node[rotate=\CF_rotatenode#1\_nil,anchor=\CF_anchornode#1\_nil#3,at=(shifted@node)]{\expandafter\CF_gobarg\CF_gobarg#1}%
			}%
		\else
			\ifboolKV[chemfig]{scheme debug}
			{%
				node[pos=#2,sloped,yshift=#3\CF_arrowlabelsep,draw,fill,cyan](shifted@node){}%
				node[draw,pos=#2,anchor=-#390,sloped,yshift=#3\CF_arrowlabelsep]{#1}%
			}
			{%
				node[pos=#2,anchor=-#390,sloped,yshift=#3\CF_arrowlabelsep]{#1}%
			}%
		\fi
		}%
}

% pose des noeuds décalés de la dimension #1 à (\CF_arrowstartnode) et (\CF_arrowendnode)
\def\CF_arrowshiftnodes#1{%
	\unless\ifdim\CF_ifempty{#1}\CF_zero{#1}=0pt
		\expanded{%
			\noexpand\path(\CF_arrowstartnode)--(\CF_arrowendnode)%
			node[pos=0,sloped,yshift=#1](\CF_arrowstartnode1){}node[pos=1,sloped,yshift=#1](\CF_arrowendnode1){};}%
		\edef\CF_arrowstartnode{\CF_arrowstartnode1}\edef\CF_arrowendnode{\CF_arrowendnode1}%
	\fi
}

\definearrow3{->}{%
	\CF_arrowshiftnodes{#3}%
	\CF_eafter{\draw[}\CF_arrowcurrentstyle](\CF_arrowstartnode)--(\CF_arrowendnode);%
	\CF_arrowdisplaylabel{#1}{0.5}+\CF_arrowstartnode{#2}{0.5}-\CF_arrowendnode
}

\definearrow3{<-}{%
	\CF_arrowshiftnodes{#3}%
	\CF_eafter{\draw[}\CF_arrowcurrentstyle](\CF_arrowendnode)--(\CF_arrowstartnode);%
	\CF_arrowdisplaylabel{#1}{0.5}+\CF_arrowstartnode{#2}{0.5}-\CF_arrowendnode
}

\definearrow5{-/>}{%
	\CF_arrowshiftnodes{#3}%
	\CF_eafter{\draw[}\CF_arrowcurrentstyle](\CF_arrowstartnode)--(\CF_arrowendnode)%
		coordinate[midway,shift=(\CF_arrowcurrentangle:-1pt)](midway@i)%
		coordinate[midway,shift=(\CF_arrowcurrentangle:1pt)](midway@ii)%
		coordinate[at=(midway@i),shift=(\CF_ifempty{#4}{225}{#4+180}+\CF_arrowcurrentangle:\CF_ifempty{#5}{5pt}{#5})](line@start)%
		coordinate[at=(midway@i),shift=(\CF_ifempty{#4}{45}{#4}+\CF_arrowcurrentangle:\CF_ifempty{#5}{5pt}{#5})](line@end)%
		coordinate[at=(midway@ii),shift=(\CF_ifempty{#4}{225}{#4+180}+\CF_arrowcurrentangle:\CF_ifempty{#5}{5pt}{#5})](line@start@i)%
		coordinate[at=(midway@ii),shift=(\CF_ifempty{#4}{45}{#4}+\CF_arrowcurrentangle:\CF_ifempty{#5}{5pt}{#5})](line@end@i);
	\draw(line@start)--(line@end);%
	\draw(line@start@i)--(line@end@i);%
	\CF_arrowdisplaylabel{#1}{0.5}+\CF_arrowstartnode{#2}{0.5}-\CF_arrowendnode
}

\definearrow3{<->}{%
	\CF_arrowshiftnodes{#3}%
	\CF_eafter{\draw[}\CF_arrowcurrentstyle,\CF_arrowtip-\CF_arrowtip](\CF_arrowstartnode)--(\CF_arrowendnode);%
	\CF_arrowdisplaylabel{#1}{0.5}+\CF_arrowstartnode{#2}{0.5}-\CF_arrowendnode
}

\definearrow3{<=>}{%
	\CF_arrowshiftnodes{#3}%
	\path[allow upside down](\CF_arrowstartnode)--(\CF_arrowendnode)%
			node[pos=0,sloped,yshift=\CF_arrowdoublesep](\CF_arrowstartnode @u0){}%
			node[pos=0,sloped,yshift=-\CF_arrowdoublesep](\CF_arrowstartnode @d0){}%
			node[pos=1,sloped,yshift=\CF_arrowdoublesep](\CF_arrowstartnode @u1){}%
			node[pos=1,sloped,yshift=-\CF_arrowdoublesep](\CF_arrowstartnode @d1){};%
	\begingroup
		\ifboolKV[chemfig]{arrow double harpoon}
			\pgfarrowharpoontrue
			{}%
		\CF_eafter{\draw[}\CF_arrowcurrentstyle](\CF_arrowstartnode @u0)--(\CF_arrowstartnode @u1);%
		\CF_eafter{\draw[}\CF_arrowcurrentstyle](\CF_arrowstartnode @d1)--(\CF_arrowstartnode @d0);%
	\endgroup
	\CF_arrowdisplaylabel{#1}{0.5}+\CF_arrowstartnode{#2}{0.5}-\CF_arrowendnode%
}

\definearrow3{<->>}{%
	\CF_arrowshiftnodes{#3}%
	\path[allow upside down](\CF_arrowstartnode)--(\CF_arrowendnode)%
			node[pos=0,sloped,yshift=1pt](\CF_arrowstartnode @u0){}%
			node[pos=\CF_arrowdoubleposstart,sloped,yshift=-1pt](\CF_arrowstartnode @d0){}%
			node[pos=1,sloped,yshift=1pt](\CF_arrowstartnode @u1){}%
			node[pos=\CF_arrowdoubleposend,sloped,yshift=-1pt](\CF_arrowstartnode @d1){};%
	\begingroup
		\ifboolKV[chemfig]{arrow double harpoon}
			\pgfarrowharpoontrue
			{}%
		\CF_eafter{\draw[}\CF_arrowcurrentstyle](\CF_arrowstartnode @u0)--(\CF_arrowstartnode @u1);%
		\CF_eafter{\draw[}\CF_arrowcurrentstyle](\CF_arrowstartnode @d1)--(\CF_arrowstartnode @d0);%
	\endgroup
	\CF_arrowdisplaylabel{#1}{0.5}+\CF_arrowstartnode{#2}{0.5}-\CF_arrowendnode%
}

\definearrow3{<<->}{%
	\path[allow upside down](\CF_arrowstartnode)--(\CF_arrowendnode)%
			node[pos=\CF_arrowdoubleposstart,sloped,yshift=1pt](\CF_arrowstartnode @u0){}%
			node[pos=0,sloped,yshift=-1pt](\CF_arrowstartnode @d0){}%
			node[pos=\CF_arrowdoubleposend,sloped,yshift=1pt](\CF_arrowstartnode @u1){}%
			node[pos=1,sloped,yshift=-1pt](\CF_arrowstartnode @d1){};%
	\begingroup
		\ifboolKV[chemfig]{arrow double harpoon}
			\pgfarrowharpoontrue
			{}%
		\CF_eafter{\draw[}\CF_arrowcurrentstyle](\CF_arrowstartnode @u0)--(\CF_arrowstartnode @u1);%
		\CF_eafter{\draw[}\CF_arrowcurrentstyle](\CF_arrowstartnode @d1)--(\CF_arrowstartnode @d0);%
	\endgroup
	\CF_arrowdisplaylabel{#1}{0.5}+\CF_arrowstartnode{#2}{0.5}-\CF_arrowendnode
}

\definearrow30{%
	\CF_arrowshiftnodes{#3}%
	\CF_arrowdisplaylabel{#1}{0.5}+\CF_arrowstartnode{#2}{0.5}-\CF_arrowendnode
}

\definearrow5{-U>}{%
	\CF_arrowshiftnodes{#3}%
	\CF_eafter{\draw[}\CF_arrowcurrentstyle](\CF_arrowstartnode)--(\CF_arrowendnode)node[midway](Uarrow@arctangent){};%
	\CF_ifempty{#4}
	{%
		\def\CF_Uarrowradius{0.333}%
	}
	{%
		\def\CF_Uarrowradius{#4}%
	}%
	\CF_ifempty{#5}%
	{%
		\def\CF_Uarrowabsangle{60}%
	}
	{%
		\pgfmathsetmacro\CF_Uarrowabsangle{abs(#5)}% ne prendre en compte que la valeur absolue de l'angle
	}%
	\expandafter\draw\expanded{[\CF_ifempty{#1}{draw=none}{\unexpanded\expandafter{\CF_arrowcurrentstyle}},-]}(Uarrow@arctangent)%
		arc[radius=\CF_compoundsep*\CF_currentarrowlength*\CF_Uarrowradius,start angle=\CF_arrowcurrentangle-90,delta angle=-\CF_Uarrowabsangle]node(Uarrow@start){};
	\expandafter\draw\expanded{[\CF_ifempty{#2}{draw=none}{\unexpanded\expandafter{\CF_arrowcurrentstyle}}]}(Uarrow@arctangent)%
		arc[radius=\CF_compoundsep*\CF_currentarrowlength*\CF_Uarrowradius,start angle=\CF_arrowcurrentangle-90,delta angle=\CF_Uarrowabsangle]node(Uarrow@end){};
	\pgfmathsetmacro\CF_temp{\CF_Uarrowradius*cos(\CF_arrowcurrentangle)<0?"-":"+"}%
	\ifdim\CF_Uarrowradius pt>0pt
		\CF_arrowdisplaylabel{#1}{0}\CF_temp{Uarrow@start}{#2}{1}\CF_temp{Uarrow@end}%
	\else
		\CF_arrowdisplaylabel{#2}{0}\CF_temp{Uarrow@start}{#1}{1}\CF_temp{Uarrow@end}%
	\fi
}

\def\CF_grabdelim#1#2#3\_CFnil{\def\CF_leftdelim{#1}\def\CF_rightdelim{#2}}
\defKV[CFdelimiters]{%
	delimiters   = \CF_grabdelim#1()\_CFnil,
	height       = \def\CF_delimheight{#1},
	depth        = \CF_esecond{\CF_defifempty\CF_delimdepth{#1}}\CF_delimheight,
	open xshift  = \edef\CF_leftdelimxshift{\the\dimexpr#1},
	close xshift = \edef\CF_rightdelimxshift{\CF_ifempty{#1}{-\CF_leftdelimxshift}{-\the\dimexpr#1}}
}
\setKVdefault[CFdelimiters]{%
	delimiters   = (),
	height       = 10pt,
	depth        = ,
	open xshift  = 0pt,
	close xshift = ,
	h align      = true,
	auto rotate  = false,
	rotate       = 0,
	indice       = n
}%
\def\polymerdelim{\CF_ifnexttok[{\CF_polymerdelima}{\CF_polymerdelima[]}}
\def\CF_polymerdelima[#1]#2#3{%
	\restoreKV[CFdelimiters]%
	\CF_doifnotempty{#1}
	{%
		\setKV[CFdelimiters]{#1}%
	}%
	\edef\CF_delimhalfdim{\the\dimexpr(\CF_delimheight+\CF_delimdepth)/2}%
	\edef\CF_delimvshift {\the\dimexpr(\CF_delimheight-\CF_delimdepth)/2}%
	\chemmove{%
		\nulldelimiterspace0pt
		\CF_extractxy{#2}{center}\CF_leftdelimx\CF_leftdelimy
		\CF_extractxy{#3}{center}\CF_rightdelimx\CF_rightdelimy
		\def\CF_autorotate{0}%
		\ifboolKV[CFdelimiters]{h align}
		{%
			\let\CF_rightdelimy\CF_leftdelimy
		}
		{%
			\ifboolKV[CFdelimiters]{auto rotate}
			{%
				\pgfmathatantwo{\CF_rightdelimy-\CF_leftdelimy}{\CF_rightdelimx-\CF_leftdelimx}%
				\let\CF_autorotate\pgfmathresult
			}
			{%
				\CF_eesecond\CF_ifempty{\useKV[CFdelimiters]{rotate}}
				{}
				{%
					\edef\CF_autorotate{\useKV[CFdelimiters]{rotate}}%
				}%
			}%
		}%
		\node[at={(\CF_leftdelimx+\CF_leftdelimxshift,\CF_leftdelimy+\CF_delimvshift)},rotate=\CF_autorotate]
			{$\left\CF_leftdelim\vrule height\CF_delimhalfdim depth\CF_delimhalfdim width0pt\right.$};%
		\node[at={(\CF_rightdelimx+\CF_rightdelimxshift,\CF_rightdelimy+\CF_delimvshift)},rotate=\CF_autorotate]
			{$\left.\vrule height\CF_delimhalfdim depth\CF_delimhalfdim width0pt\right\CF_rightdelim
				\CF_eesecond\CF_doifnotempty{\useKV[CFdelimiters]{indice}}
					{\CF_underscore{\rlap{$\scriptstyle\useKV[CFdelimiters]{indice}$}}}
			$};
	}%
}

\catcode`\@11
\pgfdeclarearrow{%
	name = CF,%
	defaults = {%
		length  = 3pt 5 1,%
		width'  = 0pt .8,%
		inset'  = 0pt .5,%
		line width = 0pt 1 1,%
		round%
	},%
	setup code = {%
		% Cap the line width at 1/4th distance from inset to tip
		\pgf@x\pgfarrowlength
		\advance\pgf@x by-\pgfarrowinset
		\pgf@x.25\pgf@x
		\ifdim\pgf@x<\pgfarrowlinewidth\pgfarrowlinewidth\pgf@x\fi
		% Compute front miter length:
		\pgfmathdivide@{\pgf@sys@tonumber\pgfarrowlength}{\pgf@sys@tonumber\pgfarrowwidth}%
		\let\pgf@temp@quot\pgfmathresult%
		\pgf@x\pgfmathresult pt%
		\pgf@x\pgfmathresult\pgf@x%
		\pgf@x4\pgf@x%
		\advance\pgf@x by1pt%
		\pgfmathsqrt@{\pgf@sys@tonumber\pgf@x}%
		\pgf@xc\pgfmathresult\pgfarrowlinewidth% xc is front miter
		\pgf@xc.5\pgf@xc
		\pgf@xa\pgf@temp@quot\pgfarrowlinewidth% xa is extra harpoon miter
		% Compute back miter length:
		\pgf@ya.5\pgfarrowwidth%
		\csname pgfmathatan2@\endcsname{\pgfmath@tonumber\pgfarrowlength}{\pgfmath@tonumber\pgf@ya}%
		\pgf@yb\pgfmathresult pt%
		\csname pgfmathatan2@\endcsname{\pgfmath@tonumber\pgfarrowinset}{\pgfmath@tonumber\pgf@ya}%
		\pgf@ya\pgfmathresult pt%
		\advance\pgf@yb by-\pgf@ya%
		\pgf@yb.5\pgf@yb% half angle in yb
		\pgfmathtan@{\pgf@sys@tonumber\pgf@yb}%
		\pgfmathreciprocal@{\pgfmathresult}%
		\pgf@yc\pgfmathresult\pgfarrowlinewidth%
		\pgf@yc.5\pgf@yc%
		\advance\pgf@ya by\pgf@yb%
		\pgfmathsincos@{\pgf@sys@tonumber\pgf@ya}%
		\pgf@ya\pgfmathresulty\pgf@yc% ya is the back miter
		\pgf@yb\pgfmathresultx\pgf@yc% yb is the top miter
		\ifdim\pgfarrowinset=0pt
			\pgf@ya.5\pgfarrowlinewidth% easy: back miter is half linewidth
		\fi
		% Compute inset miter length:
		\pgfmathdivide@{\pgf@sys@tonumber\pgfarrowinset}{\pgf@sys@tonumber\pgfarrowwidth}%
		\let\pgf@temp@quot\pgfmathresult%
		\pgf@x\pgfmathresult pt%
		\pgf@x\pgfmathresult\pgf@x%
		\pgf@x4\pgf@x%\pgf@ya
		\advance\pgf@x by1pt%
		\pgfmathsqrt@{\pgf@sys@tonumber\pgf@x}%
		\pgf@yc\pgfmathresult\pgfarrowlinewidth% yc is inset miter
		\pgf@yc.5\pgf@yc% 
		% Inner length (pgfutil@tempdima) is now arrowlength - front miter - back miter
		\pgfutil@tempdima\pgfarrowlength%
		\advance\pgfutil@tempdima by-\pgf@xc%
		\advance\pgfutil@tempdima by-\pgf@ya%
		\pgfutil@tempdimb.5\pgfarrowwidth%
		\advance\pgfutil@tempdimb by-\pgf@yb%
		% harpoon miter correction
		\ifpgfarrowroundjoin
			\pgfarrowssetbackend{\pgf@ya\advance\pgf@x by-.5\pgfarrowlinewidth}%
		\else
			\pgfarrowssetbackend{0pt}
		\fi
		\ifpgfarrowharpoon
			\pgfarrowssetlineend{\pgfarrowinset\advance\pgf@x by\pgf@yc\advance\pgf@x by.5\pgfarrowlinewidth}%
		\else
			\pgfarrowssetlineend{\pgfarrowinset\advance\pgf@x by\pgf@yc\advance\pgf@x by-.25\pgfarrowlinewidth}%
			\ifpgfarrowreversed
				\ifdim\pgfinnerlinewidth>0pt
					\pgfarrowssetlineend{\pgfarrowinset}%
				\else
					\pgfarrowssetlineend{\pgfutil@tempdima\advance\pgf@x by\pgf@ya\advance\pgf@x by-.25\pgfarrowlinewidth}%
				\fi
			\fi
		\fi
		\ifpgfarrowroundjoin
			\pgfarrowssettipend{\pgfutil@tempdima\advance\pgf@x by\pgf@ya\advance\pgf@x by.5\pgfarrowlinewidth}%
		\else
			\pgfarrowssettipend{\pgfarrowlength\ifpgfarrowharpoon\advance\pgf@x by\pgf@xa\fi}%
		\fi
		% The hull:
		\pgfarrowshullpoint{\pgfarrowlength\ifpgfarrowroundjoin\else\ifpgfarrowharpoon\advance\pgf@x by\pgf@xa\fi\fi}{\ifpgfarrowharpoon-.5\pgfarrowlinewidth\else0pt\fi}%
		\pgfarrowsupperhullpoint{0pt}{.5\pgfarrowwidth}%
		\pgfarrowshullpoint{\pgfarrowinset}{\ifpgfarrowharpoon-.5\pgfarrowlinewidth\else 0pt\fi}%
		% Adjust inset
		\pgfarrowssetvisualbackend{\pgfarrowinset}%
		\advance\pgfarrowinset by\pgf@yc%
		% The following are needed in the code:
		\pgfarrowssavethe\pgfutil@tempdima
		\pgfarrowssavethe\pgfutil@tempdimb
		\pgfarrowssavethe\pgfarrowlinewidth
		\pgfarrowssavethe\pgf@ya
		\pgfarrowssavethe\pgfarrowinset
	},%
	drawing code = {%
		\pgfsetdash{}{0pt}%
		\ifpgfarrowroundjoin\pgfsetroundjoin\else\pgfsetmiterjoin\fi
		\ifdim\pgfarrowlinewidth=\pgflinewidth\else\pgfsetlinewidth{\pgfarrowlinewidth}\fi
		\pgfpathmoveto{\pgfqpoint{\pgfutil@tempdima\advance\pgf@x by\pgf@ya}{0pt}}%
		\pgfpathlineto{\pgfqpoint{\pgf@ya}{\pgfutil@tempdimb}}%
		\pgfpathlineto{\pgfqpoint{\pgfarrowinset}{0pt}}%
		\ifpgfarrowharpoon \else
			\pgfpathlineto{\pgfqpoint{\pgf@ya}{-\pgfutil@tempdimb}}%
		\fi
		\pgfpathclose
		\ifpgfarrowopen\pgfusepathqstroke\else\ifdim\pgfarrowlinewidth>0pt \pgfusepathqfillstroke\else\pgfusepathqfill\fi\fi
	},%
	parameters = {%
		\the\pgfarrowlinewidth,%
		\the\pgfarrowlength,%
		\the\pgfarrowwidth,%
		\the\pgfarrowinset,%
		\ifpgfarrowharpoon h\fi%
		\ifpgfarrowopen o\fi%
		\ifpgfarrowroundjoin j\fi%
	}%
}
\CF_restorecatcode
\endinput